\documentclass[UTF8, 12pt, fontset=fandol, oneside]{ctexart}
\usepackage{researchnotes}

\title{第七章\quad 相似标准型}
\author{}
\date{}

\begin{document}

\maketitle

\section*{§7.1\quad 多项式矩阵}

我们将在这一章里继续探讨上一章中提出的问题：给定一个线性变换，找出一组基，使该线性变换在这组基下的表示矩阵具有比较简单的形状. 这个问题等价于寻找一类比较简单的矩阵，使任一同阶方阵均与这类矩阵中的某一个相似. 这类比较简单的矩阵就是所谓的相似标准型.

为了解决这个问题，可以分两步走. 第一步找出相似矩阵的不变量，这些不变量不仅在相似关系下保持不变，而且足以判断两个矩阵是否相似. 我们称这样的不变量为全系不变量. 比如秩是两个同阶矩阵在相抵关系下的不变量，反之，若两个同阶矩阵的秩相同，则它们必相抵. 因此，秩是矩阵相抵关系的全系不变量. 第二步找出一类比较简单的矩阵，利用相似关系的全系不变量就可以判断一个矩阵与这类矩阵中的某一个相似.

相似关系比相抵关系要更复杂一些，它的全系不变量也比较复杂. 我们在上一章中已经知道，矩阵的特征多项式（从而特征值）是相似不变量，但它并不是全系不变量，因为我们很容易举出例子来证明这一点. 比如下面两个矩阵的特征多项式相同但不相似：
\[
A=\left(\begin{array}{cc}
0 & 1\\
0 & 0
\end{array}\right),\quad
B=\left(\begin{array}{cc}
0 & 0\\
0 & 0
\end{array}\right).
\]
$A$ 的特征多项式与 $B$ 的特征多项式都是 $\lambda^2$，但 $A$ 与 $B$ 绝不相似.

人们经过研究终于发现，两个矩阵 $A$ 与 $B$ 之间的相似和 $\lambda I_n-A$ 与 $\lambda I_n-B$ 的相抵有着密切的联系. 注意，$\lambda I_n-A$ 是这样形式的矩阵：
\[
\left(\begin{array}{cccc}
\lambda-a_{11} & -a_{12} & \cdots & -a_{1n}\\
-a_{21} & \lambda-a_{22} & \cdots & -a_{2n}\\
\vdots & \vdots & & \vdots\\
-a_{n1} & -a_{n2} & \cdots & \lambda-a_{nn}
\end{array}\right),
\]
其中的元素含有未定元 $\lambda$. 一般地，下列形式的矩阵：
\[
A(\lambda)=
\left(\begin{array}{cccc}
a_{11}(\lambda) & a_{12}(\lambda) & \cdots & a_{1n}(\lambda)\\
a_{21}(\lambda) & a_{22}(\lambda) & \cdots & a_{2n}(\lambda)\\
\vdots & \vdots & & \vdots\\
a_{m1}(\lambda) & a_{m2}(\lambda) & \cdots & a_{mn}(\lambda)
\end{array}\right),
\]
其中 $a_{ij}(\lambda)$ 是以 $\lambda$ 为未定元的数域 $\mathbb{K}$ 上的多项式，称为多项式矩阵，或 $\lambda$-矩阵. $\lambda$-矩阵的加法、数乘及乘法与数域上的矩阵运算一样，只需在运算过程中将数的运算代之以多项式运算即可.

现在我们来研究两个 $\lambda$-矩阵的相抵关系. 首先我们必须定义什么叫 $\lambda$-矩阵的初等变换.

\noindent\textbf{定义 7.1.1}\quad 对 $\lambda$-矩阵 $A(\lambda)$ 施行的下列 3 种变换称为 $\lambda$-矩阵的初等行变换：

(1) 将 $A(\lambda)$ 的两行对换；

(2) 将 $A(\lambda)$ 的第 $i$ 行乘以 $\mathbb{K}$ 中的非零常数 $c$；

(3) 将 $A(\lambda)$ 的第 $i$ 行乘以 $\mathbb{K}$ 上的多项式 $f(\lambda)$ 后加到第 $j$ 行上去.

同理我们可以定义 3 种 $\lambda$-矩阵的初等列变换.

\noindent\textbf{定义 7.1.2}\quad 若 $A(\lambda),B(\lambda)$ 是同阶 $\lambda$-矩阵且 $A(\lambda)$ 经过 $\lambda$-矩阵的初等变换后可变为 $B(\lambda)$，则称 $A(\lambda)$ 与 $B(\lambda)$ 相抵.

与数字矩阵一样，$\lambda$-矩阵的相抵关系也是一种等价关系，即

(1) $A(\lambda)$ 与自身相抵；

(2) 若 $A(\lambda)$ 与 $B(\lambda)$ 相抵，则 $B(\lambda)$ 与 $A(\lambda)$ 相抵；

(3) 若 $A(\lambda)$ 与 $B(\lambda)$ 相抵，$B(\lambda)$ 与 $C(\lambda)$ 相抵，则 $A(\lambda)$ 与 $C(\lambda)$ 相抵.

证明与数域上相同，请读者自己完成.

类似数字矩阵，$\lambda$-矩阵的初等变换也对应于初等 $\lambda$-矩阵的相乘.

\noindent\textbf{定义 7.1.3}\quad 下列 3 种矩阵称为初等 $\lambda$-矩阵：

(1) 将 $n$ 阶单位阵的第 $i$ 行与第 $j$ 行对换，记为 $P_{ij}$；

(2) 将 $n$ 阶单位阵的第 $i$ 行乘以非零常数 $c$，记为 $P_i(c)$；

(3) 将 $n$ 阶单位阵的第 $i$ 行乘以多项式 $f(\lambda)$ 后加到第 $j$ 行上去得到的矩阵，记为 $T_{ij}(f(\lambda))$.

注意，第一类与第二类初等 $\lambda$-矩阵与数域上的第一类与第二类初等矩阵没有什么区别. 第三类初等 $\lambda$-矩阵的形状如下：
\[
T_{ij}(f(\lambda))=
\left(\begin{array}{cccccc}
1 & & & & & \\
& \ddots & & & & \\
& & 1 & & & \\
& & \vdots & \ddots & & \\
& & f(\lambda) & \cdots & 1 & \\
& & & & & \ddots\\
& & & & & 1
\end{array}\right).
\]

\noindent\textbf{定理 7.1.1}\quad 对 $\lambda$-矩阵 $A(\lambda)$ 施行第 $k$ ($k=1,2,3$) 类初等行（列）变换等于用第 $k$ 类初等 $\lambda$-矩阵左（右）乘以 $A(\lambda)$.

\noindent\textbf{证明}\quad 与定理 2.4.3 完全相同，留给读者作为练习.

\noindent\textbf{注}\quad 下列 $\lambda$-矩阵的变换不是 $\lambda$-矩阵的初等变换：
\[
\left(\begin{array}{cc}
1 & 1\\
0 & 1
\end{array}\right)\to
\left(\begin{array}{cc}
\lambda & \lambda\\
0 & 1
\end{array}\right).
\]
这是因为前面一个矩阵的第一行乘以 $\lambda$ 不是 $\lambda$-矩阵的初等变换. 同理下面的变换需第一行乘以 $\lambda^{-1}$，因此也不是 $\lambda$-矩阵的初等变换：
\[
\left(\begin{array}{cc}
\lambda & 0\\
0 & 1
\end{array}\right)\to
\left(\begin{array}{cc}
1 & 0\\
0 & 1
\end{array}\right).
\]

对 $n$ 阶 $\lambda$-矩阵，我们可定义可逆 $\lambda$-矩阵的概念.

\noindent\textbf{定义 7.1.4}\quad 若 $A(\lambda),B(\lambda)$ 都是 $n$ 阶 $\lambda$-矩阵，且
\[
A(\lambda)B(\lambda)=B(\lambda)A(\lambda)=I_n,
\]
则称 $B(\lambda)$ 是 $A(\lambda)$ 的逆 $\lambda$-矩阵. 这时称 $A(\lambda)$ 为可逆 $\lambda$-矩阵，在不引起混淆的情形下，有时简称为可逆阵.

\noindent\textbf{注}\quad 注意不要将数字矩阵中的一些结论随意搬到 $\lambda$-矩阵上. 比如下面的 $\lambda$-矩阵的行列式不为零，但它不是可逆 $\lambda$-矩阵：
\[
\left(\begin{array}{cc}
\lambda & 0\\
0 & 1
\end{array}\right).
\]
这是因为矩阵
\[
\left(\begin{array}{cc}
\lambda^{-1} & 0\\
0 & 1
\end{array}\right)
\]
不是 $\lambda$-矩阵之故.

容易证明，有限个可逆 $\lambda$-矩阵之积仍是可逆 $\lambda$-矩阵，而初等 $\lambda$-矩阵都是可逆 $\lambda$-矩阵，因此有限个初等 $\lambda$-矩阵之积也是可逆 $\lambda$-矩阵. 下一节我们将证明可逆 $\lambda$-矩阵必可表示为有限个初等 $\lambda$-矩阵之积.

为了把数域上矩阵的相似关系归结为 $\lambda$-矩阵的相抵关系，我们尚需证明一个有关 $\lambda$-矩阵带余除法的引理. 设 $M(\lambda)$ 是一个 $n$ 阶 $\lambda$-矩阵，则 $M(\lambda)$ 可以化为如下形状：
\[
M(\lambda)=M_m\lambda^m+M_{m-1}\lambda^{m-1}+\cdots+M_0,
\]
其中 $M_i$ 为数域 $\mathbb{K}$ 上的 $n$ 阶数字矩阵. 因此，一个多项式矩阵可以化为系数为矩阵的多项式，反之亦然.

\noindent\textbf{引理 7.1.1}\quad 设 $M(\lambda)$ 与 $N(\lambda)$ 是两个 $n$ 阶 $\lambda$-矩阵且都不等于零，又设 $B$ 为 $n$ 阶数字矩阵，则必存在 $\lambda$-矩阵 $Q(\lambda)$ 及 $S(\lambda)$ 和数字矩阵 $R$ 及 $T$，使下式成立：
\begin{equation}
M(\lambda)=(\lambda I-B)Q(\lambda)+R,
\tag{7.1.1}
\end{equation}
\begin{equation}
N(\lambda)=S(\lambda)(\lambda I-B)+T.
\tag{7.1.2}
\end{equation}

\noindent\textbf{证明}\quad 将 $M(\lambda)$ 写为
\[
M(\lambda)=M_m\lambda^m+M_{m-1}\lambda^{m-1}+\cdots+M_0,
\]
其中 $M_m\neq O$. 可对 $m$ 用归纳法，若 $m=0$，则已适合要求（取 $Q(\lambda)=O$）. 现设对小于 $m$ 次的矩阵多项式，(7.1.1) 式成立. 令
\[
Q_1(\lambda)=M_m\lambda^{m-1},
\]
则
\begin{equation}
M(\lambda)-(\lambda I-B)Q_1(\lambda)=(BM_m+M_{m-1})\lambda^{m-1}+\cdots+M_0.
\tag{7.1.3}
\end{equation}
上式是一个次数小于 $m$ 的矩阵多项式，由归纳假设得
\[
M(\lambda)-(\lambda I-B)Q_1(\lambda)=(\lambda I-B)Q_2(\lambda)+R.
\]
于是
\[
M(\lambda)=(\lambda I-B)\left[Q_1(\lambda)+Q_2(\lambda)\right]+R.
\]
令 $Q(\lambda)=Q_1(\lambda)+Q_2(\lambda)$ 即得 (7.1.1) 式. 同理可证 (7.1.2) 式. $\square$

\noindent\textbf{定理 7.1.2}\quad 设 $A,B$ 是数域 $\mathbb{K}$ 上的矩阵，则 $A$ 与 $B$ 相似的充分必要条件是 $\lambda$-矩阵 $\lambda I-A$ 与 $\lambda I-B$ 相抵.

\noindent\textbf{证明}\quad 若 $A$ 与 $B$ 相似，则存在 $\mathbb{K}$ 上的非异阵 $P$，使 $B=P^{-1}AP$，于是
\begin{equation}
P^{-1}(\lambda I-A)P=\lambda I-P^{-1}AP=\lambda I-B.
\tag{7.1.4}
\end{equation}
把 $P$ 看成是常数 $\lambda$-矩阵，上式表明 $\lambda I-A$ 与 $\lambda I-B$ 相抵.

反过来，若 $\lambda I-A$ 与 $\lambda I-B$ 相抵，即存在 $M(\lambda)$ 及 $N(\lambda)$，使
\begin{equation}
M(\lambda)(\lambda I-A)N(\lambda)=\lambda I-B,
\tag{7.1.5}
\end{equation}
其中 $M(\lambda)$ 与 $N(\lambda)$ 都是有限个初等矩阵之积，因而都是可逆阵. 因此可将 (7.1.5) 式写为
\begin{equation}
M(\lambda)(\lambda I-A)=(\lambda I-B)N(\lambda)^{-1},
\tag{7.1.6}
\end{equation}
由引理 7.1.1 可设
\begin{equation}
M(\lambda)=(\lambda I-B)Q(\lambda)+R,
\tag{7.1.7}
\end{equation}
代入 (7.1.6) 式经整理得
\begin{equation}
R(\lambda I-A)=(\lambda I-B)\left[N(\lambda)^{-1}-Q(\lambda)(\lambda I-A)\right].
\tag{7.1.8}
\end{equation}
上式左边是次数小于等于 1 的矩阵多项式，因此上式右边中括号内的矩阵多项式的次数必须小于等于零，也即必是一个常数矩阵，设为 $P$. 于是
\begin{equation}
R(\lambda I-A)=(\lambda I-B)P.
\tag{7.1.9}
\end{equation}
(7.1.9) 式又可整理为
\[
(R-P)\lambda=RA-BP.
\]
再次比较次数得 $R=P,RA=BP$. 现只需证明 $P$ 是一个非异阵即可. 由假设，
\[
P=N(\lambda)^{-1}-Q(\lambda)(\lambda I-A),
\]
将上式两边右乘 $N(\lambda)$ 并移项得
\[
PN(\lambda)+Q(\lambda)(\lambda I-A)N(\lambda)=I.
\]

但
\[
(\lambda I-A)N(\lambda)=M(\lambda)^{-1}(\lambda I-B),
\]
因此
\begin{equation}
PN(\lambda)+Q(\lambda)M(\lambda)^{-1}(\lambda I-B)=I.
\tag{7.1.10}
\end{equation}
再由引理 7.1.1 可设
\[
N(\lambda)=S(\lambda)(\lambda I-B)+T,
\]
代入 (7.1.10) 式并整理得
\[
\left[PS(\lambda)+Q(\lambda)M(\lambda)^{-1}\right](\lambda I-B)=I-PT.
\]
上式右边是次数小于等于零的矩阵多项式，因此上式左边中括号内的矩阵多项式必须为零，从而 $PT=I$，即 $P$ 是非异阵. $\square$

\subsection*{习题 7.1}

1. 设 $M(\lambda)$ 是 $\lambda$-矩阵且可写为
\[
M(\lambda)=M_m\lambda^m+M_{m-1}\lambda^{m-1}+\cdots+M_0.
\]
求证：若 $M(\lambda)$ 是可逆 $\lambda$-矩阵，则 $M_0$ 是非异阵.

2. 证明引理 7.1.1 中的余式 $R$ 及 $T$ 是唯一确定的，$Q(\lambda)$ 与 $S(\lambda)$ 也唯一确定.

3. 设 $A,B$ 是数域 $\mathbb{K}$ 上的 $n$ 阶矩阵，求证：它们的特征矩阵 $\lambda I_n-A$ 和 $\lambda I_n-B$ 相抵的充分必要条件是存在同阶矩阵 $P$ 和 $Q$，使 $A=PQ$，$B=QP$，且 $P$ 及 $Q$ 中至少有一个是可逆阵.

4. 设 $A,B$ 是数域 $\mathbb{K}$ 上的 $n$ 阶矩阵，$\lambda I_n-A$ 相抵于 $\operatorname{diag}\{f_1(\lambda),f_2(\lambda),\cdots,f_n(\lambda)\}$，$\lambda I_n-B$ 相抵于 $\operatorname{diag}\{f_{i_1}(\lambda),f_{i_2}(\lambda),\cdots,f_{i_n}(\lambda)\}$，其中 $f_{i_1}(\lambda),f_{i_2}(\lambda),\cdots,f_{i_n}(\lambda)$ 是 $f_1(\lambda),f_2(\lambda),\cdots,f_n(\lambda)$ 的一个排列. 求证：$A$ 与 $B$ 相似.

\section*{§7.2\quad 矩阵的法式}

在上一节中，我们把矩阵的相似归结为 $\lambda$-矩阵的相抵. 现在我们来求 $\lambda$-矩阵的相抵标准型，我们自然地希望任一 $\lambda$-矩阵相抵于一个对角 $\lambda$-矩阵.

\noindent\textbf{引理 7.2.1}\quad 设 $A(\lambda)=(a_{ij}(\lambda))_{m\times n}$ 是任一非零 $\lambda$-矩阵，则 $A(\lambda)$ 必相抵于这样的一个 $\lambda$-矩阵 $B(\lambda)=(b_{ij}(\lambda))_{m\times n}$，其中 $b_{11}(\lambda)\neq 0$ 且 $b_{11}(\lambda)$ 可整除 $B(\lambda)$ 中的任一元素 $b_{ij}(\lambda)$.

\noindent\textbf{证明}\quad 设 $k=\min\{\deg a_{ij}(\lambda)\mid a_{ij}(\lambda)\neq 0,1\leq i\leq m;1\leq j\leq n\}$，我们对 $k$ 用数学归纳法. 首先，经行对换及列对换可将 $A(\lambda)$ 的第 $(1,1)$ 元素变成次数最低的非零多项式，因此不妨设 $a_{11}(\lambda)\neq 0$ 且 $\deg a_{11}(\lambda)=k$. 若 $k=0$，则 $a_{11}(\lambda)$ 是一个非零常数，结论显然成立. 假设对非零元素次数的最小值小于 $k$ 的任一 $\lambda$-矩阵，引理的结论成立，现考虑非零元素次数的最小值等于 $k$ 的 $\lambda$-矩阵 $A(\lambda)$. 若 $a_{11}(\lambda)$ 可整除所有的 $a_{ij}(\lambda)$，则结论已成立. 若否，设在第一列中有元素 $a_{i1}(\lambda)$ 不能被 $a_{11}(\lambda)$ 整除，作带余除法：
\[
a_{i1}(\lambda)=a_{11}(\lambda)q(\lambda)+r(\lambda).
\]
用 $-q(\lambda)$ 乘以第一行加到第 $i$ 行上，第 $(i,1)$ 元素就变为 $r(\lambda)$. 注意到 $r(\lambda)\neq 0$ 且 $\deg r(\lambda)<\deg a_{11}(\lambda)=k$，由归纳假设即知结论成立.

同样的方法可施于第一行. 因此我们不妨设 $a_{11}(\lambda)$ 可整除第一行及第一列. 这时，设 $a_{21}(\lambda)=a_{11}(\lambda)q(\lambda)$. 将第一行乘以 $-q(\lambda)$ 加到第二行上，则第 $(2,1)$ 元素变为零. 用同样的方法可消去第一行、第一列除 $a_{11}(\lambda)$ 以外的所有元素，于是 $A(\lambda)$ 经初等变换后变成下列形状：
\begin{equation}
\left(\begin{array}{cccc}
a_{11}(\lambda) & 0 & \cdots & 0\\
0 & a'_{22}(\lambda) & \cdots & a'_{2n}(\lambda)\\
\vdots & \vdots & & \vdots\\
0 & a'_{m2}(\lambda) & \cdots & a'_{mn}(\lambda)
\end{array}\right).
\tag{7.2.1}
\end{equation}
这时，若 $a_{11}(\lambda)$ 可整除所有其他元素，则结论已成立. 若否，比如 $a_{11}(\lambda)$ 不能整除 $a'_{ij}(\lambda)$，则将第 $i$ 行加到第一行上去，这时在第一行又出现了一元素 $a'_{ij}(\lambda)$，它不能被 $a_{11}(\lambda)$ 整除. 重复上面的做法，通过归纳假设即可得到结论. $\square$

\noindent\textbf{定理 7.2.1}\quad 设 $A(\lambda)$ 是一个 $n$ 阶 $\lambda$-矩阵，则 $A(\lambda)$ 相抵于对角阵
\begin{equation}
\operatorname{diag}\{d_1(\lambda),d_2(\lambda),\cdots,d_r(\lambda);0,\cdots,0\},
\tag{7.2.2}
\end{equation}
其中 $d_i(\lambda)$ 是非零首一多项式且 $d_i(\lambda)\mid d_{i+1}(\lambda)$ ($i=1,2,\cdots,r-1$).

\noindent\textbf{证明}\quad 对 $n$ 用数学归纳法，当 $n=1$ 时结论显然，现设 $A(\lambda)$ 是 $n$ 阶 $\lambda$-矩阵. 由引理可知 $A(\lambda)$ 相抵于 $n$ 阶 $\lambda$-矩阵 $B(\lambda)=(b_{ij}(\lambda))$，其中 $b_{11}(\lambda)\mid b_{ij}(\lambda)$ 对一切 $i,j$ 成立. 因此，将 $B(\lambda)$ 的第一行乘以 $\lambda$ 的某个多项式加到第二行上去便可消去 $b_{21}(\lambda)$. 同理可依次消去第一列除 $b_{11}(\lambda)$ 以外的所有元素. 再用类似方法消去第一行其余元素. 这样便得到了一个矩阵：
\[
\left(\begin{array}{cccc}
b_{11}(\lambda) & 0 & \cdots & 0\\
0 & b'_{22}(\lambda) & \cdots & b'_{2n}(\lambda)\\
\vdots & \vdots & & \vdots\\
0 & b'_{n2}(\lambda) & \cdots & b'_{nn}(\lambda)
\end{array}\right).
\]
不难看出，这时 $b_{11}(\lambda)$ 仍可整除所有的 $b'_{ij}(\lambda)$. 设 $c$ 为 $b_{11}(\lambda)$ 的首项系数，记 $d_1(\lambda)=c^{-1}b_{11}(\lambda)$，设 $\overline{B}(\lambda)$ 为上面的矩阵中右下方的 $n-1$ 阶 $\lambda$-矩阵，则由归纳假设可知存在 $P(\lambda)$ 及 $Q(\lambda)$，使
\[
P(\lambda)\overline{B}(\lambda)Q(\lambda)=\operatorname{diag}\{d_2(\lambda),\cdots,d_r(\lambda);0,\cdots,0\},
\]
且 $d_i(\lambda)\mid d_{i+1}(\lambda)$ ($i=2,\cdots,r-1$)，其中 $P(\lambda)$ 与 $Q(\lambda)$ 可写成为有限个 $n-1$ 阶初等 $\lambda$-矩阵之积. 因此
\[
\left(\begin{array}{cc}
1 & O\\
O & P(\lambda)
\end{array}\right)
\left(\begin{array}{cc}
d_1(\lambda) & O\\
O & \overline{B}(\lambda)
\end{array}\right)
\left(\begin{array}{cc}
1 & O\\
O & Q(\lambda)
\end{array}\right)
=\operatorname{diag}\{d_1(\lambda),d_2(\lambda),\cdots,d_r(\lambda);0,\cdots,0\},
\]
且
\[
\left(\begin{array}{cc}
1 & O\\
O & P(\lambda)
\end{array}\right),\quad
\left(\begin{array}{cc}
1 & O\\
O & Q(\lambda)
\end{array}\right)
\]
可写成有限个 $n$ 阶初等 $\lambda$-矩阵之积. 于是只需证明 $d_1(\lambda)\mid d_2(\lambda)$ 即可. 但这点很容易看出，事实上由于 $\overline{B}(\lambda)$ 中的任一元素均可被 $d_1(\lambda)$ 整除，因此 $P(\lambda)\overline{B}(\lambda)Q(\lambda)$ 中的任一元素也可被 $d_1(\lambda)$ 整除，这就证明了定理. $\square$

\noindent\textbf{注}\quad 我们上面对 $n$ 阶 $\lambda$-矩阵证明了它必相抵于一个对角阵. 事实上，对长方 $\lambda$-矩阵，结论也同样成立，证明也类似. (7.2.2) 式中的 $r$ 通常称为 $A(\lambda)$ 的秩. 但要注意即使某个 $n$ 阶 $\lambda$-矩阵的秩等于 $n$，它也未必是可逆 $\lambda$-矩阵.

\noindent\textbf{推论 7.2.1}\quad 任一 $n$ 阶可逆 $\lambda$-矩阵都可表示为有限个初等 $\lambda$-矩阵之积.

\noindent\textbf{证明}\quad 由上述定理，存在 $P(\lambda),Q(\lambda)$，使可逆阵 $A(\lambda)$ 适合
\[
P(\lambda)A(\lambda)Q(\lambda)=\operatorname{diag}\{d_1(\lambda),d_2(\lambda),\cdots,d_r(\lambda);0,\cdots,0\},
\]
其中 $P(\lambda),Q(\lambda)$ 为有限个初等 $\lambda$-矩阵之积. 因为上式左边是个可逆阵，故右边的矩阵也可逆，从而 $r=n$. 注意一个对角 $\lambda$-矩阵要可逆必须 $d_1(\lambda),d_2(\lambda),\cdots,d_n(\lambda)$ 皆为非零常数，又它们都是首一多项式，故只能是 $d_1(\lambda)=d_2(\lambda)=\cdots=d_n(\lambda)=1$，于是
\[
A(\lambda)=P(\lambda)^{-1}Q(\lambda)^{-1}.
\]
因为初等 $\lambda$-矩阵的逆仍是初等 $\lambda$-矩阵，故 $P(\lambda)^{-1}$ 与 $Q(\lambda)^{-1}$ 都是有限个初等 $\lambda$-矩阵之积，从而 $A(\lambda)$ 也由有限个初等 $\lambda$-矩阵之积. $\square$

\noindent\textbf{推论 7.2.2}\quad 设 $A$ 是数域 $\mathbb{K}$ 上的 $n$ 阶矩阵，则 $A$ 的特征矩阵 $\lambda I_n-A$ 必相抵于
\begin{equation}
\operatorname{diag}\{1,\cdots,1,d_1(\lambda),\cdots,d_m(\lambda)\},
\tag{7.2.3}
\end{equation}
其中 $d_i(\lambda)\mid d_{i+1}(\lambda)$ ($i=1,2,\cdots,m-1$).

\noindent\textbf{证明}\quad 由上述定理，存在 $P(\lambda),Q(\lambda)$，使
\[
P(\lambda)(\lambda I_n-A)Q(\lambda)=\operatorname{diag}\{d_1(\lambda),d_2(\lambda),\cdots,d_r(\lambda);0,\cdots,0\},
\]
其中 $P(\lambda),Q(\lambda)$ 为有限个初等 $\lambda$-矩阵之积，根据 $\lambda$-矩阵初等变换的定义以及行列式的性质可得，上式左边的行列式等于 $c|\lambda I_n-A|$，其中 $c$ 是一个非零常数，从而上式右边的行列式不为零，故 $r=n$. 把 $d_i(\lambda)$ 中的常数多项式写出来（因是首一多项式，故为常数 1），即得结论. $\square$

\noindent\textbf{定义 7.2.1}\quad 称 (7.2.2) 式中的对角 $\lambda$-矩阵为 $A(\lambda)$ 的法式或相抵标准型.

\noindent\textbf{例 7.2.1}\quad 求 $\lambda I-A$ 的法式，其中
\[
A=\left(\begin{array}{ccc}
0 & 1 & -1\\
3 & -2 & 0\\
-1 & 1 & -1
\end{array}\right).
\]

\noindent\textbf{解}\quad
\[
\lambda I-A=
\left(\begin{array}{ccc}
\lambda & -1 & 1\\
-3 & \lambda+2 & 0\\
1 & -1 & \lambda+1
\end{array}\right)
\sim
\left(\begin{array}{ccc}
1 & -1 & \lambda+1\\
-3 & \lambda+2 & 0\\
\lambda & -1 & 1
\end{array}\right)
\]
\[
\sim
\left(\begin{array}{ccc}
1 & -1 & \lambda+1\\
0 & \lambda-1 & 3\lambda+3\\
0 & \lambda-1 & -\lambda^2-\lambda+1
\end{array}\right)
\sim
\left(\begin{array}{ccc}
1 & 0 & 0\\
0 & \lambda-1 & 3\lambda+3\\
0 & \lambda-1 & -\lambda^2-\lambda+1
\end{array}\right)
\]
\[
\sim
\left(\begin{array}{ccc}
1 & 0 & 0\\
0 & \lambda-1 & 6\\
0 & \lambda-1 & -\lambda^2-4\lambda+4
\end{array}\right)
\sim
\left(\begin{array}{ccc}
1 & 0 & 0\\
0 & 6 & \lambda-1\\
0 & -\lambda^2-4\lambda+4 & \lambda-1
\end{array}\right)
\]
\[
\sim
\left(\begin{array}{ccc}
1 & 0 & 0\\
0 & 6 & 6(\lambda-1)\\
0 & -\lambda^2-4\lambda+4 & 6(\lambda-1)
\end{array}\right)
\sim
\left(\begin{array}{ccc}
1 & 0 & 0\\
0 & 6 & 0\\
0 & -\lambda^2-4\lambda+4 & (\lambda-1)(\lambda^2+4\lambda+2)
\end{array}\right)
\]
\[
\sim
\left(\begin{array}{ccc}
1 & 0 & 0\\
0 & 1 & 0\\
0 & -\lambda^2-4\lambda+4 & (\lambda-1)(\lambda^2+4\lambda+2)
\end{array}\right)
\sim
\left(\begin{array}{ccc}
1 & 0 & 0\\
0 & 1 & 0\\
0 & 0 & (\lambda-1)(\lambda^2+4\lambda+2)
\end{array}\right).
\]

\subsection*{习题 7.2}

1. 试求 $\lambda I-A$ 的法式，其中 $A$ 为
\[
(1)\ A=\left(\begin{array}{ccc}
1 & 2 & -1\\
0 & 1 & 1\\
0 & 0 & -2
\end{array}\right);\quad
(2)\ A=\left(\begin{array}{ccc}
-3 & 1 & 1\\
-3 & 0 & 2\\
-2 & 1 & 0
\end{array}\right);\quad
(3)\ A=\left(\begin{array}{ccc}
0 & 1 & 1\\
0 & 1 & 0\\
-1 & 1 & 2
\end{array}\right).
\]

2. 用初等变换化 $\lambda$-矩阵为对角阵且适合定理 7.2.1 的要求：
\[
(1)\ A(\lambda)=\left(\begin{array}{ccc}
1-\lambda & 2\lambda-1 & \lambda\\
\lambda & \lambda^2 & -\lambda\\
1+\lambda^2 & \lambda^2+\lambda-1 & -\lambda^2
\end{array}\right);
\]
\[
(2)\ A(\lambda)=\left(\begin{array}{cccc}
0 & 0 & 0 & \lambda^2\\
0 & 0 & \lambda^2-\lambda & 0\\
0 & (\lambda-1)^2 & 0 & 0\\
\lambda^2-\lambda & 0 & 0 & 0
\end{array}\right).
\]

3. 设 $f(\lambda),g(\lambda)$ 是数域 $\mathbb{K}$ 上的首一多项式，$d(\lambda)=(f(\lambda),g(\lambda))$，$m(\lambda)=[f(\lambda),g(\lambda)]$ 分别是 $f(\lambda),g(\lambda)$ 的最大公因式和最小公倍式，证明下列 3 个 $\lambda$-矩阵相抵：
\[
\left(\begin{array}{cc}
f(\lambda) & 0\\
0 & g(\lambda)
\end{array}\right),\quad
\left(\begin{array}{cc}
g(\lambda) & 0\\
0 & f(\lambda)
\end{array}\right),\quad
\left(\begin{array}{cc}
d(\lambda) & 0\\
0 & m(\lambda)
\end{array}\right).
\]

4. 求证：$n$ 阶 $\lambda$-矩阵为可逆阵的充分必要条件是它的行列式为非零常数.

5. 设 $\lambda I-A$ 的法式为 $\operatorname{diag}\{1,\cdots,1,d_1(\lambda),\cdots,d_m(\lambda)\}$，求证：$A$ 的特征多项式
\[
|\lambda I-A|=d_1(\lambda)\cdots d_m(\lambda).
\]

\section*{§7.3\quad 不变因子}

在上一节中我们证明了任一 $\lambda$-矩阵均相抵于一对角 $\lambda$-矩阵. 因此，如果两个 $n$ 阶 $\lambda$-矩阵的法式相同，则它们必相抵. 现在要问反过来的问题，即如果两个 $\lambda$-矩阵的法式不相同，是否它们必不相抵？假如我们能证明这一点，那么我们就找到了 $\lambda$-矩阵相抵关系的全系不变量，即 $r$ 个首一多项式序列：
\begin{equation}
d_1(\lambda),d_2(\lambda),\cdots,d_r(\lambda)
\tag{7.3.1}
\end{equation}
适合 $d_i(\lambda)\mid d_{i+1}(\lambda)$ ($i=1,\cdots,r-1$). 为了证明这一点，我们只要证明 (7.3.1) 式中的多项式在相抵关系下具有不变性就可以了. 为此，我们需引进行列式因子的概念.

\noindent\textbf{定义 7.3.1}\quad 设 $A(\lambda)$ 是 $n$ 阶 $\lambda$-矩阵，$k$ 是小于等于 $n$ 的正整数. 如果 $A(\lambda)$ 有一个 $k$ 阶子式不为零，则定义 $A(\lambda)$ 的 $k$ 阶行列式因子 $D_k(\lambda)$ 为 $A(\lambda)$ 的所有 $k$ 阶子式的最大公因式（首一多项式）. 如果 $A(\lambda)$ 的所有 $k$ 阶子式都等于零，则定义 $A(\lambda)$ 的 $k$ 阶行列式因子 $D_k(\lambda)$ 为零.

\noindent\textbf{例 7.3.1}\quad 求下列矩阵的行列式因子：
\[
A(\lambda)=
\left(\begin{array}{ccccc}
d_1(\lambda) & & & & \\
& \ddots & & & \\
& & d_r(\lambda) & & \\
& & & 0 & \\
& & & & \ddots
\end{array}\right),
\]
其中 $d_i(\lambda)$ 为非零首一多项式且 $d_i(\lambda)\mid d_{i+1}(\lambda)$ ($i=1,2,\cdots,r-1$).

\noindent\textbf{解}\quad $A(\lambda)$ 的非零行列式因子为
\[
D_1(\lambda)=d_1(\lambda),\quad D_2(\lambda)=d_1(\lambda)d_2(\lambda),\quad \cdots,\quad D_r(\lambda)=d_1(\lambda)d_2(\lambda)\cdots d_r(\lambda).
\]

\noindent\textbf{引理 7.3.1}\quad 设 $D_1(\lambda),D_2(\lambda),\cdots,D_r(\lambda)$ 是 $A(\lambda)$ 的非零行列式因子，则
\[
D_i(\lambda)\mid D_{i+1}(\lambda),\quad i=1,2,\cdots,r-1.
\]

\noindent\textbf{证明}\quad 设 $A_{i+1}$ 是 $A(\lambda)$ 的任一 $i+1$ 阶子式，即在 $A(\lambda)$ 中任意取出 $i+1$ 行及 $i+1$ 列组成的行列式. 将这个行列式按某一行展开，则它的每一个展开项都是一个多项式与一个 $i$ 阶子式的乘积. 由于 $D_i(\lambda)$ 是所有 $i$ 阶子式的公因子，因此 $D_i(\lambda)\mid A_{i+1}$. 而 $D_{i+1}(\lambda)$ 是所有 $i+1$ 阶子式的最大公因子，因此 $D_i(\lambda)\mid D_{i+1}(\lambda)$ 对一切 $i=1,2,\cdots,r-1$ 成立. $\square$

\noindent\textbf{定义 7.3.2}\quad 设 $D_1(\lambda),D_2(\lambda),\cdots,D_r(\lambda)$ 是 $\lambda$-矩阵 $A(\lambda)$ 的非零行列式因子，则 $g_1(\lambda)=D_1(\lambda)$，$g_2(\lambda)=D_2(\lambda)/D_1(\lambda),\cdots,g_r(\lambda)=D_r(\lambda)/D_{r-1}(\lambda)$ 称为 $A(\lambda)$ 的不变因子.

例 7.3.1 中矩阵的不变因子为
\[
d_1(\lambda),d_2(\lambda),\cdots,d_r(\lambda).
\]

\noindent\textbf{定理 7.3.1}\quad 相抵的 $\lambda$-矩阵有相同的行列式因子，从而有相同的不变因子.

\noindent\textbf{证明}\quad 我们只需证明行列式因子在三类初等变换下不改变就可以了. 对第一类初等变换，交换 $\lambda$-矩阵 $A(\lambda)$ 的任意两行（列），显然 $A(\lambda)$ 的 $i$ 阶子式最多改变一个符号，因此行列式因子不改变.

对第二类初等变换，$A(\lambda)$ 的 $i$ 阶子式与变换后矩阵的 $i$ 阶子式最多差一个非零常数，因此行列式因子也不改变.

对第三类初等变换，记变换后的矩阵为 $B(\lambda)$，则 $B(\lambda)$ 与 $A(\lambda)$ 的 $i$ 阶子式可能出现以下 3 种情形：子式完全相同；$B(\lambda)$ 子式中的某一行（列）等于 $A(\lambda)$ 中相应子式的同一行（列）加上该子式中某一行（列）与某个多项式之积；$B(\lambda)$ 子式中的某一行（列）等于 $A(\lambda)$ 中相应子式的同一行（列）加上不在该子式中的某一行（列）与某个多项式之积. 在前面两种情形，行列式的值不改变，因此不影响行列式因子. 现在来讨论第三种情形. 设 $B_i$ 为 $B(\lambda)$ 的 $i$ 阶子式，相应的 $A(\lambda)$ 的 $i$ 阶子式记为 $A_i$，则由行列式的性质得
\begin{equation}
B_i=A_i+f(\lambda)\widetilde{A}_i,
\tag{7.3.2}
\end{equation}
其中 $\widetilde{A}_i$ 由 $A(\lambda)$ 中的 $i$ 行与 $i$ 列组成，因此它与 $A(\lambda)$ 的某个 $i$ 阶子式最多差一个符号，$f(\lambda)$ 是乘以某一行（列）的那个多项式，于是 $A(\lambda)$ 的行列式因子 $D_i(\lambda)\mid A_i,D_i(\lambda)\mid \widetilde{A}_i$，故 $D_i(\lambda)\mid B_i$. 这说明，$D_i(\lambda)$ 可整除 $B(\lambda)$ 的所有 $i$ 阶子式，因此 $D_i(\lambda)$ 可整除 $B(\lambda)$ 的 $i$ 阶行列式因子 $\widetilde{D}_i(\lambda)$. 但 $B(\lambda)$ 也可用第三类初等变换变成 $A(\lambda)$，于是 $\widetilde{D}_i(\lambda)\mid D_i(\lambda)$. 由于 $D_i(\lambda)$ 及 $\widetilde{D}_i(\lambda)$ 都是首一多项式，因此必有 $D_i(\lambda)=\widetilde{D}_i(\lambda)$. $\square$

\noindent\textbf{推论 7.3.1}\quad 设 $n$ 阶 $\lambda$-矩阵 $A(\lambda)$ 的法式为
\[
\Lambda=\operatorname{diag}\{d_1(\lambda),d_2(\lambda),\cdots,d_r(\lambda);0,\cdots,0\},
\]
其中 $d_i(\lambda)$ 是非零首一多项式且 $d_i(\lambda)\mid d_{i+1}(\lambda)$ ($i=1,2,\cdots,r-1$)，则 $A(\lambda)$ 的不变因子为 $d_1(\lambda),d_2(\lambda),\cdots,d_r(\lambda)$. 特别，法式和不变因子之间相互唯一确定.

\noindent\textbf{证明}\quad 由定理 7.3.1，$A(\lambda)$ 与 $\Lambda$ 有相同的不变因子. 再由例 7.3.1，$\Lambda$ 的不变因子为 $d_1(\lambda),d_2(\lambda),\cdots,d_r(\lambda)$，从而它们也是 $A(\lambda)$ 的不变因子. $\square$

\noindent\textbf{推论 7.3.2}\quad 设 $A(\lambda),B(\lambda)$ 为 $n$ 阶 $\lambda$-矩阵，则 $A(\lambda)$ 与 $B(\lambda)$ 相抵当且仅当它们有相同的法式.

\noindent\textbf{证明}\quad 若 $A(\lambda)$ 与 $B(\lambda)$ 有相同的法式，显然它们相抵. 若 $A(\lambda)$ 与 $B(\lambda)$ 相抵，由定理 7.3.1 知 $A(\lambda)$ 与 $B(\lambda)$ 有相同的不变因子，从而有相同的法式. $\square$

\noindent\textbf{推论 7.3.3}\quad $n$ 阶 $\lambda$-矩阵 $A(\lambda)$ 的法式与初等变换的选取无关.

\noindent\textbf{证明}\quad 设 $\Lambda_1,\Lambda_2$ 是 $A(\lambda)$ 通过不同的初等变换得到的两个法式，则 $\Lambda_1$ 与 $\Lambda_2$ 相抵，由推论 7.3.2 可得 $\Lambda_1=\Lambda_2$. $\square$

\noindent\textbf{定理 7.3.2}\quad 数域 $\mathbb{K}$ 上 $n$ 阶矩阵 $A$ 与 $B$ 相似的充分必要条件是它们的特征矩阵 $\lambda I-A$ 与 $\lambda I-B$ 具有相同的行列式因子或不变因子.

\noindent\textbf{证明}\quad 显然不变因子与行列式因子之间相互唯一确定. 再由定理 7.1.2、推论 7.3.1 及推论 7.3.2 即得结论. $\square$

以后特征矩阵 $\lambda I-A$ 的行列式因子和不变因子均简称为 $A$ 的行列式因子和不变因子.

\noindent\textbf{推论 7.3.4}\quad 设 $\mathbb{F}\subseteq\mathbb{K}$ 是两个数域，$A,B$ 是 $\mathbb{F}$ 上的两个矩阵，则 $A$ 与 $B$ 在 $\mathbb{F}$ 上相似的充分必要条件是它们在 $\mathbb{K}$ 上相似.

\noindent\textbf{证明}\quad 若 $A$ 与 $B$ 在 $\mathbb{F}$ 上相似，由于 $\mathbb{F}\subseteq\mathbb{K}$，它们当然在 $\mathbb{K}$ 上也相似. 反之，若 $A$ 与 $B$ 在 $\mathbb{K}$ 上相似，则 $\lambda I-A$ 与 $\lambda I-B$ 在 $\mathbb{K}$ 上有相同的不变因子，也就是说它们有相同的法式. 由推论 7.3.3 可知，求法式与初等变换的选取无关. 注意到 $\lambda I-A$ 与 $\lambda I-B$ 是数域 $\mathbb{F}$ 上的 $\lambda$-矩阵，故可用 $\mathbb{F}$ 上 $\lambda$-矩阵的初等变换就能将它们变成法式，其中只涉及 $\mathbb{F}$ 中数的加、减、乘、除运算以及 $\mathbb{F}$ 上的多项式的加、减、乘、数乘运算，最后得到法式中的不变因子 $d_i(\lambda)$ 仍是 $\mathbb{F}$ 上的多项式. 这就是说存在 $\mathbb{F}$ 上的可逆 $\lambda$-矩阵 $P(\lambda),Q(\lambda),M(\lambda),N(\lambda)$，使
\[
P(\lambda)(\lambda I-A)Q(\lambda)=M(\lambda)(\lambda I-B)N(\lambda)=\operatorname{diag}\{d_1(\lambda),\cdots,d_n(\lambda)\},
\]
从而
\[
M(\lambda)^{-1}P(\lambda)(\lambda I-A)Q(\lambda)N(\lambda)^{-1}=\lambda I-B,
\]
即 $\lambda I-A$ 与 $\lambda I-B$ 在 $\mathbb{F}$ 上相抵，由定理 7.1.2 可得 $A$ 与 $B$ 在 $\mathbb{F}$ 上相似. $\square$

推论 7.3.4 告诉我们：矩阵的相似关系在基域扩张下不变. 事实上，推论 7.3.4 的证明过程也说明：矩阵的不变因子在基域扩张下也不变.

\subsection*{习题 7.3}

1. 求下列矩阵的行列式因子与不变因子：
\[
(1)\ \left(\begin{array}{ccc}
1 & 2 & 0\\
0 & 2 & 0\\
-2 & -2 & 1
\end{array}\right);\quad
(2)\ \left(\begin{array}{ccc}
4 & 5 & -2\\
-2 & -2 & 1\\
-1 & -1 & 1
\end{array}\right);\quad
(3)\ \left(\begin{array}{ccc}
3 & 1 & -1\\
2 & 2 & -1\\
2 & 2 & 0
\end{array}\right).
\]

2. 判断下列矩阵是否相似：
\[
(1)\ \left(\begin{array}{ccc}
1 & 1 & 0\\
0 & 1 & 0\\
0 & 0 & -1
\end{array}\right),\quad
\left(\begin{array}{ccc}
1 & 0 & 0\\
0 & 0 & 1\\
0 & 1 & 0
\end{array}\right);
\]
\[
(2)\ \left(\begin{array}{ccc}
3 & 2 & 2\\
2 & 3 & 2\\
-4 & -4 & -3
\end{array}\right),\quad
\left(\begin{array}{ccc}
0 & 1 & 1\\
0 & 1 & 0\\
-1 & 1 & 2
\end{array}\right);
\]
\[
(3)\ \left(\begin{array}{ccc}
2 & -1 & 2\\
5 & -3 & 3\\
-1 & 0 & -2
\end{array}\right),\quad
\left(\begin{array}{ccc}
-3 & 1 & 1\\
-3 & 0 & 2\\
-2 & 1 & 0
\end{array}\right).
\]

3. 求下列 $n$ 阶上三角阵的行列式因子和不变因子：
\[
\left(\begin{array}{ccccc}
a & 1 & 1 & \cdots & 1\\
0 & a & 1 & \cdots & 1\\
0 & 0 & a & \cdots & 1\\
\vdots & \vdots & \vdots & & \vdots\\
0 & 0 & 0 & \cdots & a
\end{array}\right).
\]

4. 证明：任一 $n$ 阶方阵 $A$ 必与它的转置 $A'$ 相似.

5. 设 $A$ 是 $n$ 阶方阵，证明以下三个结论等价：

(1) $A=cI_n$，其中 $c$ 为常数；

(2) $A$ 的 $n-1$ 阶行列式因子是一个 $n-1$ 次多项式；

(3) $A$ 的不变因子组中无常数.

\section*{§7.4\quad 有理标准型}

利用矩阵的不变因子，现在可以来构造所谓的“有理标准型”了. 我们的想法是寻找一个比较简单的矩阵，使它与给定的矩阵有相同的不变因子. 由前面两节我们已经知道，矩阵 $A$ 的特征矩阵 $\lambda I-A$ 的法式为
\[
\operatorname{diag}\{1,\cdots,1,d_1(\lambda),\cdots,d_k(\lambda)\},
\]
其中 $d_i(\lambda)$ 为非常数首一多项式且 $d_i(\lambda)\mid d_{i+1}(\lambda)$ ($i=1,2,\cdots,k-1$)，则 $A$ 的不变因子就是
\[
1,\cdots,1,d_1(\lambda),\cdots,d_k(\lambda).
\]

\noindent\textbf{引理 7.4.1}\quad 设 $r$ 阶矩阵
\[
F=\left(\begin{array}{ccccc}
0 & 1 & 0 & \cdots & 0\\
0 & 0 & 1 & \cdots & 0\\
\vdots & \vdots & \vdots & & \vdots\\
0 & 0 & 0 & \cdots & 1\\
-a_r & -a_{r-1} & -a_{r-2} & \cdots & -a_1
\end{array}\right),
\]
则

(1) $F$ 的行列式因子为
\begin{equation}
1,\cdots,1,f(\lambda),
\tag{7.4.1}
\end{equation}
其中共有 $r-1$ 个 $1$，$f(\lambda)=\lambda^r+a_1\lambda^{r-1}+\cdots+a_r$，$F$ 的不变因子也由 (7.4.1) 式给出；

(2) $F$ 的极小多项式等于 $f(\lambda)$.

\noindent\textbf{证明}\quad (1) $F$ 的 $r$ 阶行列式因子就是它的特征多项式，由例 1.5.7 得
\[
|\lambda I-F|=\lambda^r+a_1\lambda^{r-1}+\cdots+a_r.
\]
对任一 $k<r$，$\lambda I-F$ 总有一个 $k$ 阶子式其值等于 $(-1)^k$，故 $D_k(\lambda)=1$.

(2) 因为 $F$ 的特征多项式为 $f(\lambda)$，所以 $F$ 适合多项式 $f(\lambda)$. 设 $e_i$ ($i=1,2,\cdots,r$) 是 $r$ 维标准单位行向量，则不难算出：
\[
e_1F=e_2,\quad e_1F^2=e_3,\quad \cdots,\quad e_1F^{r-1}=e_r.
\]
显然，$e_1,e_1F,\cdots,e_1F^{r-1}$ 是一组线性无关的向量，因此 $F$ 不可能适合一个次数不超过 $r-1$ 的非零多项式，从而 $F$ 的极小多项式就是 $f(\lambda)$. $\square$

\noindent\textbf{引理 7.4.2}\quad 设 $\lambda$-矩阵 $A(\lambda)$ 相抵于对角 $\lambda$-矩阵
\begin{equation}
\operatorname{diag}\{d_1(\lambda),d_2(\lambda),\cdots,d_n(\lambda)\},
\tag{7.4.2}
\end{equation}
$\lambda$-矩阵 $B(\lambda)$ 相抵于对角 $\lambda$-矩阵
\begin{equation}
\operatorname{diag}\{d'_1(\lambda),d'_2(\lambda),\cdots,d'_n(\lambda)\},
\tag{7.4.3}
\end{equation}
且 $d'_1(\lambda),d'_2(\lambda),\cdots,d'_n(\lambda)$ 是 $d_1(\lambda),d_2(\lambda),\cdots,d_n(\lambda)$ 的一个置换（即若不计次序，这两组多项式完全相同），则 $A(\lambda)$ 相抵于 $B(\lambda)$.

\noindent\textbf{证明}\quad 利用行对换及列对换即可将 (7.4.2) 式变成 (7.4.3) 式，因此 (7.4.2) 式所示的矩阵与 (7.4.3) 式所示的矩阵相抵，从而 $A(\lambda)$ 与 $B(\lambda)$ 相抵. $\square$

\noindent\textbf{定理 7.4.1}\quad 设 $A$ 是数域 $\mathbb{K}$ 上的 $n$ 阶方阵，$A$ 的不变因子组为
\[
1,\cdots,1,d_1(\lambda),\cdots,d_k(\lambda),
\]
其中 $\deg d_i(\lambda)=m_i\geq 1$，则 $A$ 相似于下列分块对角阵：
\begin{equation}
F=\left(\begin{array}{cccc}
F_1 & & & \\
& F_2 & & \\
& & \ddots & \\
& & & F_k
\end{array}\right),
\tag{7.4.4}
\end{equation}
其中 $F_i$ 的阶等于 $m_i$，且 $F_i$ 是形如引理 7.4.1 中的矩阵，$F_i$ 的最后一行由 $d_i(\lambda)$ 的系数（除首项系数之外）的负值组成.

\noindent\textbf{证明}\quad 注意到 $\lambda I-A$ 的第 $n$ 个行列式因子就是 $A$ 的特征多项式 $|\lambda I-A|$，再由行列式因子的相抵不变性可知：
\[
|\lambda I-A|=d_1(\lambda)d_2(\lambda)\cdots d_k(\lambda).
\]
而 $|\lambda I-A|$ 是一个 $n$ 次多项式，因此 $m_1+m_2+\cdots+m_k=n$. 一方面，$\lambda I-A$ 的法式为
\[
\operatorname{diag}\{1,\cdots,1,d_1(\lambda),d_2(\lambda),\cdots,d_k(\lambda)\},
\]
其中有 $n-k$ 个 $1$. 另一方面，对 $\lambda I-F$ 的每个分块都施以 $\lambda$-矩阵的初等变换，由引理 7.4.1 可知，$\lambda I-F$ 相抵于如下对角阵：
\begin{equation}
\operatorname{diag}\{1,\cdots,1,d_1(\lambda);1,\cdots,1,d_2(\lambda);\cdots;1,\cdots,1,d_k(\lambda)\},
\tag{7.4.5}
\end{equation}
其中每个 $d_i(\lambda)$ 前各有 $m_i-1$ 个 $1$，从而共有 $\sum_{i=1}^k(m_i-1)=n-k$ 个 $1$. 因此 (7.4.5) 式所示的矩阵与 $\lambda I-A$ 的法式只相差主对角线上元素的置换，由引理 7.4.2 可得 $\lambda I-A$ 与 $\lambda I-F$ 相抵，从而 $A$ 与 $F$ 相似. $\square$

\noindent\textbf{定义 7.4.1}\quad (7.4.4) 式称为矩阵 $A$ 的有理标准型或 Frobenius 标准型，每个 $F_i$ 称为 Frobenius 块.

\noindent\textbf{例 7.4.1}\quad 设 6 阶矩阵 $A$ 的不变因子为
\[
1,1,1,\lambda-1,(\lambda-1)^2,(\lambda-1)^2(\lambda+1),
\]
则 $A$ 的有理标准型为
\[
\left(\begin{array}{cccccc}
1 & & & & & \\
& 0 & 1 & & & \\
& -1 & 2 & & & \\
& & & 0 & 1 & 0\\
& & & 0 & 0 & 1\\
& & & -1 & 1 & 1
\end{array}\right).
\]

在第六章第 3 节中我们定义了一个矩阵的极小多项式，用有理标准型可以清楚地表明极小多项式与不变因子的关系.

\noindent\textbf{定理 7.4.2}\quad 设数域 $\mathbb{K}$ 上的 $n$ 阶矩阵 $A$ 的不变因子为
\[
1,\cdots,1,d_1(\lambda),\cdots,d_k(\lambda),
\]
其中 $d_i(\lambda)\mid d_{i+1}(\lambda)$ ($i=1,\cdots,k-1$)，则 $A$ 的极小多项式 $m(\lambda)=d_k(\lambda)$.

\noindent\textbf{证明}\quad 设 $A$ 的有理标准型为
\[
F=\left(\begin{array}{cccc}
F_1 & & & \\
& F_2 & & \\
& & \ddots & \\
& & & F_k
\end{array}\right).
\]
因为相似矩阵有相同的极小多项式，故只需证明 $F$ 的极小多项式是 $d_k(\lambda)$ 即可. 但 $F$ 是分块对角阵，由命题 6.3.3 知 $F$ 的极小多项式是诸 $F_i$ 极小多项式的最小公倍式. 又由引理 7.4.1 知 $F_i$ 的极小多项式为 $d_i(\lambda)$. 因为 $d_i(\lambda)\mid d_{i+1}(\lambda)$，故诸 $d_i(\lambda)$ 的最小公倍式等于 $d_k(\lambda)$. $\square$

\noindent\textbf{例 7.4.2}\quad 下面两个 4 阶矩阵
\[
A=\left(\begin{array}{cccc}
0 & 0 & 0 & 0\\
0 & 0 & 0 & 0\\
0 & 0 & 0 & 1\\
0 & 0 & 0 & 0
\end{array}\right),\quad
B=\left(\begin{array}{cccc}
0 & 1 & 0 & 0\\
0 & 0 & 0 & 0\\
0 & 0 & 0 & 1\\
0 & 0 & 0 & 0
\end{array}\right)
\]
的不变因子分别为 $A:1,\lambda,\lambda,\lambda^2$ 和 $B:1,1,\lambda^2,\lambda^2$. 它们的特征多项式和极小多项式分别相等，但它们不相似.

\subsection*{习题 7.4}

1. 根据下列不变因子组写出有理标准型：

(1) $1,1,\lambda,\lambda(\lambda+1)^2$；

(2) $1,1,1,(\lambda+1),(\lambda+1)^2,(\lambda+1)^3$.

2. 求下列矩阵的有理标准型：
\[
(1)\ \left(\begin{array}{ccc}
1 & 1 & 1\\
0 & 2 & 0\\
-1 & 1 & 3
\end{array}\right);\quad
(2)\ \left(\begin{array}{ccc}
1 & 3 & 3\\
3 & 1 & 3\\
-3 & -3 & -5
\end{array}\right);\quad
(3)\ \left(\begin{array}{ccc}
3 & -1 & -1\\
7 & -2 & -3\\
3 & -1 & -1
\end{array}\right).
\]

3. 若将有理标准型 $F=\operatorname{diag}\{F_1,F_2,\cdots,F_k\}$ 中任意两个 Frobenius 块 $F_i$ 与 $F_j$ 互换位置，问：所得的矩阵与 $F$ 是否相似？

4. 例 7.4.2 说明特征多项式和极小多项式分别相等的两个矩阵可能不相似. 若矩阵的阶数不超过 3，则特征多项式和极小多项式分别相等的两个矩阵是否仍有可能不相似？

5. 证明：$n$ 阶矩阵 $A$ 的特征多项式和极小多项式相等的充分必要条件是 $\lambda I_n-A$ 的行列式因子为 $1,\cdots,1,D_n(\lambda)$.

6. 设矩阵 $A$ 的特征多项式和极小多项式相等，矩阵 $B$ 也具有这个性质，又 $A$ 和 $B$ 的特征多项式相同，问：$A,B$ 是否必相似？

7. 若 $k$ 是满足 $A^k=O$ 的最小正整数，则称 $A$ 为 $k$ 次幂零阵，请写出 $A$ 的最后一个不变因子，并证明：所有 $n$ 阶 $n$ 次幂零阵彼此相似.

8. 若 $n$ 阶矩阵 $A$ 有 $n$ 个不同的特征值，求证：$A$ 的特征多项式与极小多项式相等.

9. 设数域 $\mathbb{K}$ 上的 $n$ 阶矩阵 $A$ 的特征多项式 $f(\lambda)=P_1(\lambda)P_2(\lambda)\cdots P_k(\lambda)$，其中 $P_i(\lambda)$ ($i=1,\cdots,k$) 是 $\mathbb{K}$ 上互异的首一不可约多项式. 求证：$A$ 的有理标准型只有一个 Frobenius 块，并且 $A$ 在复数域上可对角化.

10. 设数域 $\mathbb{K}$ 上的 $n$ 阶矩阵 $A$ 的不变因子是 $1,\cdots,1,d_1(\lambda),\cdots,d_k(\lambda)$，其中 $d_i(\lambda)$ 是非常数首一多项式，$d_i(\lambda)\mid d_{i+1}(\lambda)$ ($i=1,\cdots,k-1$). 求证：对 $A$ 的任一特征值 $\lambda_0$，
\[
\operatorname{r}(\lambda_0I_n-A)=n-\sum_{i=1}^k\delta_{d_i(\lambda_0),0},
\]
其中记号 $\delta_{a,b}$ 表示：若 $a=b$，取值为 1；若 $a\neq b$，取值为 0.

11. 设 $n$ 阶矩阵 $A$ 的不变因子为 $d_1(\lambda),d_2(\lambda),\cdots,d_n(\lambda)$，其中 $d_i(\lambda)\mid d_{i+1}(\lambda)$ ($i=1,\cdots,n-1$)，又 $\lambda_0$ 是 $A$ 的特征值. 求证：$\operatorname{r}(\lambda_0I_n-A)=r$ 的充分必要条件是 $(\lambda-\lambda_0)\nmid d_r(\lambda)$ 但 $(\lambda-\lambda_0)\mid d_{r+1}(\lambda)$.

12. 设 $\varphi$ 是数域 $\mathbb{K}$ 上 $n$ 维线性空间 $V$ 上的线性变换，其极小多项式的次数等于 $n$，又 $\psi$ 是 $V$ 上的另一个线性变换，满足 $\psi\varphi=\varphi\psi$. 求证：$\psi=g(\varphi)$，其中 $g(x)$ 是一个 $\mathbb{K}$ 上次数不超过 $n-1$ 的多项式.

\section*{§7.5\quad 初等因子}

利用矩阵的不变因子，我们可以求一个矩阵的有理标准型. 有理标准型对任何数域 $\mathbb{K}$ 都可以求出来，它有着诸多用途. 但是有理标准型也有一些缺点，主要是它有时不够“简单”，即有时每个 Frobenius 块太大，用起来不太方便. 有理标准型中 Frobenius 块太大的原因是不变因子 $d_i(\lambda)$ 的次数可能比较高. 如果我们用因式分解的方法分解每个 $d_i(\lambda)$，这就有可能造出更“细”的标准型来. 为此，我们先引进初等因子的概念.

设 $d_1(\lambda),d_2(\lambda),\cdots,d_k(\lambda)$ 是数域 $\mathbb{K}$ 上矩阵 $A$ 的非常数不变因子，在 $\mathbb{K}$ 上把 $d_i(\lambda)$ 分解成不可约因式之积：
\begin{equation}
\begin{aligned}
d_1(\lambda)&=p_1(\lambda)^{e_{11}}p_2(\lambda)^{e_{12}}\cdots p_t(\lambda)^{e_{1t}},\\
d_2(\lambda)&=p_1(\lambda)^{e_{21}}p_2(\lambda)^{e_{22}}\cdots p_t(\lambda)^{e_{2t}},\\
&\cdots\cdots\cdots\\
d_k(\lambda)&=p_1(\lambda)^{e_{k1}}p_2(\lambda)^{e_{k2}}\cdots p_t(\lambda)^{e_{kt}},
\end{aligned}
\tag{7.5.1}
\end{equation}
其中 $e_{ij}$ 是非负整数（注意 $e_{ij}$ 可以为零）. 由于 $d_i(\lambda)\mid d_{i+1}(\lambda)$，因此
\[
e_{1j}\leq e_{2j}\leq\cdots\leq e_{kj},\quad j=1,2,\cdots,t.
\]

\noindent\textbf{定义 7.5.1}\quad 若 (7.5.1) 式中的 $e_{ij}>0$，则称 $p_j(\lambda)^{e_{ij}}$ 为 $A$ 的一个初等因子，$A$ 的全体初等因子称为 $A$ 的初等因子组.

由因式分解的唯一性可知 $A$ 的初等因子被 $A$ 的不变因子唯一确定. 反过来，若给定一组初等因子 $p_j(\lambda)^{e_{ij}}$，适当增加一些 $1$（表示为 $p_j(\lambda)^{e_{ij}}$，其中 $e_{ij}=0$），则可将这组初等因子按不可约因式的降幂排列如下：
\begin{equation}
\begin{array}{cccc}
p_1(\lambda)^{e_{k1}}, & p_1(\lambda)^{e_{k-1,1}}, & \cdots, & p_1(\lambda)^{e_{11}},\\
p_2(\lambda)^{e_{k2}}, & p_2(\lambda)^{e_{k-1,2}}, & \cdots, & p_2(\lambda)^{e_{12}},\\
\cdots\cdots\cdots\\
p_t(\lambda)^{e_{kt}}, & p_t(\lambda)^{e_{k-1,t}}, & \cdots, & p_t(\lambda)^{e_{1t}},
\end{array}
\tag{7.5.2}
\end{equation}
令
\[
\begin{aligned}
d_k(\lambda)&=p_1(\lambda)^{e_{k1}}p_2(\lambda)^{e_{k2}}\cdots p_t(\lambda)^{e_{kt}},\\
d_{k-1}(\lambda)&=p_1(\lambda)^{e_{k-1,1}}p_2(\lambda)^{e_{k-1,2}}\cdots p_t(\lambda)^{e_{k-1,t}},\\
&\cdots\cdots\cdots\\
d_1(\lambda)&=p_1(\lambda)^{e_{11}}p_2(\lambda)^{e_{12}}\cdots p_t(\lambda)^{e_{1t}},
\end{aligned}
\]
则 $d_i(\lambda)\mid d_{i+1}(\lambda)$ ($i=1,\cdots,k-1$)，且 $d_1(\lambda),\cdots,d_k(\lambda)$ 的初等因子组就如 (7.5.2) 式所示. 因此，给定 $A$ 的初等因子组，我们可唯一地确定 $A$ 的不变因子组. 这一事实表明，$A$ 的不变因子组与初等因子组在讨论矩阵相似关系中的作用是相同的. 因此我们有下述定理.

\noindent\textbf{定理 7.5.1}\quad 数域 $\mathbb{K}$ 上的两个矩阵 $A$ 与 $B$ 相似的充分必要条件是它们有相同的初等因子组，即矩阵的初等因子组是矩阵相似关系的全系不变量.

\noindent\textbf{例 7.5.1}\quad 设 9 阶矩阵 $A$ 的不变因子组为
\[
1,\cdots,1,(\lambda-1)(\lambda^2+1),(\lambda-1)^2(\lambda^2+1)(\lambda^2-2),
\]
试分别在有理数域、实数域和复数域上求 $A$ 的初等因子组.

\noindent\textbf{解}\quad $A$ 在有理数域上的初等因子组为
\[
\lambda-1,(\lambda-1)^2,\lambda^2+1,\lambda^2+1,\lambda^2-2.
\]
$A$ 在实数域上的初等因子组为
\[
\lambda-1,(\lambda-1)^2,\lambda^2+1,\lambda^2+1,\lambda+\sqrt{2},\lambda-\sqrt{2}.
\]
$A$ 在复数域上的初等因子组为
\[
\lambda-1,(\lambda-1)^2,\lambda+\mathrm{i},\lambda+\mathrm{i},\lambda-\mathrm{i},\lambda-\mathrm{i},\lambda+\sqrt{2},\lambda-\sqrt{2}.
\]

\noindent\textbf{例 7.5.2}\quad 设 $A$ 是一个 10 阶矩阵，它的初等因子组为
\[
\lambda-1,\lambda-1,(\lambda-1)^2,(\lambda+1)^2,(\lambda+1)^3,\lambda-2.
\]
求 $A$ 的不变因子组.

\noindent\textbf{解}\quad 将上述多项式按不可约因式的降幂排列：
\[
\begin{array}{ccc}
(\lambda-1)^2, & \lambda-1, & \lambda-1,\\
(\lambda+1)^3, & (\lambda+1)^2, & 1,\\
\lambda-2, & 1, & 1.
\end{array}
\]
于是
\[
d_3(\lambda)=(\lambda-1)^2(\lambda+1)^3(\lambda-2),\quad
d_2(\lambda)=(\lambda-1)(\lambda+1)^2,\quad
d_1(\lambda)=\lambda-1.
\]
从而 $A$ 的不变因子组为
\[
1,\cdots,1,\lambda-1,(\lambda-1)(\lambda+1)^2,(\lambda-1)^2(\lambda+1)^3(\lambda-2),
\]
其中有 7 个 $1$.

用初等因子组我们可以得到比有理标准型更精细的标准型. 对每个初等因子 $p(\lambda)^l$，可构造一个比较简单的矩阵，使它的初等因子组就是 $p(\lambda)^l$，再将所有这样的矩阵拼成一个分块对角阵就可以得到标准型. 显而易见，数域越“大”，则矩阵的初等因子越多，从而分块也越精细. 我们在这里不打算构造一般数域上以初等因子为基础的标准型，在下一节中我们将讨论复数域上以初等因子为基础的 Jordan（若当）标准型.

\subsection*{习题 7.5}

1. 已知矩阵的不变因子组，求它的初等因子组（分别在有理数域、实数域与复数域上）：

(1) $1,\cdots,1,\lambda,\lambda^2(\lambda-1),\lambda^2(\lambda-1)^3(\lambda+1)$；

(2) $1,\cdots,1,\lambda^2+1,\lambda(\lambda^2+1)^2,\lambda^2(\lambda^4-1)^2$；

(3) $1,\cdots,1,\lambda(\lambda^2+1),\lambda^2(\lambda^2+1),\lambda^3(\lambda^2+1)^2(\lambda^2-2)$.

2. 已知矩阵的初等因子组，求它的不变因子组：

(1) $\lambda,\lambda,\lambda^2,\lambda+1,(\lambda+1)^2,\lambda-1,\lambda-1,(\lambda-1)^2$；

(2) $\lambda,\lambda^2,\lambda^3,\lambda-\sqrt{2},(\lambda-\sqrt{2})^2,\lambda+\sqrt{2},(\lambda+\sqrt{2})^2$；

(3) $\lambda-1,(\lambda-1)^3,\lambda+1,\lambda+1,(\lambda+1)^3,\lambda-2,(\lambda-2)^2$.

\section*{§7.6\quad Jordan 标准型}

我们根据上一节末提出的寻找标准型的思想来讨论复数域上的标准型. 由于任一多项式在复数域上均可分解为一次因子的乘积，因此复数域上的初等因子都是一次因子的幂. 又因为初等因子必是矩阵特征多项式的因式，故必具有 $(\lambda-\lambda_0)^r$ 的形状，其中 $\lambda_0$ 是矩阵的特征值. 我们先来找一个形状比较简单的矩阵，它的初等因子组就是 $(\lambda-\lambda_0)^r$.

\noindent\textbf{引理 7.6.1}\quad $r$ 阶矩阵
\[
J=\left(\begin{array}{ccccc}
\lambda_0 & 1 & & & \\
& \lambda_0 & 1 & & \\
& & \ddots & \ddots & \\
& & & \ddots & 1\\
& & & & \lambda_0
\end{array}\right)
\]
的初等因子组为 $(\lambda-\lambda_0)^r$.

\noindent\textbf{证明}\quad 显然 $J$ 的特征多项式为 $(\lambda-\lambda_0)^r$. 对任一小于 $r$ 的正整数 $k$，$\lambda I-J$ 总有一个 $k$ 阶子式，其值等于 $(-1)^k$，因此 $J$ 的行列式因子为
\begin{equation}
1,\cdots,1,(\lambda-\lambda_0)^r.
\tag{7.6.1}
\end{equation}
(7.6.1) 式也是 $J$ 的不变因子组，故 $J$ 的初等因子组只有一个多项式 $(\lambda-\lambda_0)^r$. $\square$

\noindent\textbf{引理 7.6.2}\quad 设特征矩阵 $\lambda I-A$ 经过初等变换化为下列对角阵：
\begin{equation}
\left(\begin{array}{cccc}
f_1(\lambda) & & & \\
& f_2(\lambda) & & \\
& & \ddots & \\
& & & f_n(\lambda)
\end{array}\right),
\tag{7.6.2}
\end{equation}
其中 $f_i(\lambda)$ ($i=1,\cdots,n$) 为非零首一多项式. 将 $f_i(\lambda)$ 作不可约分解，若 $(\lambda-\lambda_0)^k$ 能整除 $f_i(\lambda)$，但 $(\lambda-\lambda_0)^{k+1}$ 不能整除 $f_i(\lambda)$，就称 $(\lambda-\lambda_0)^k$ 是 $f_i(\lambda)$ 的一个准素因子，则矩阵 $A$ 的初等因子组等于所有 $f_i(\lambda)$ 的准素因子组.

\noindent\textbf{证明}\quad 第一步，先证明下列事实：

若 $f_i(\lambda),f_j(\lambda)$ ($i\neq j$) 的最大公因式和最小公倍式分别为 $g(\lambda),h(\lambda)$，则
\[
\operatorname{diag}\{f_1(\lambda),\cdots,f_i(\lambda),\cdots,f_j(\lambda),\cdots,f_n(\lambda)\}
\]
经过初等变换可以变为
\[
\operatorname{diag}\{f_1(\lambda),\cdots,g(\lambda),\cdots,h(\lambda),\cdots,f_n(\lambda)\},
\]
且这两个对角阵具有相同的准素因子组.

不失一般性，令 $i=1,j=2$. 因为 $(f_1(\lambda),f_2(\lambda))=g(\lambda)$，所以存在 $u(\lambda),v(\lambda)$，使
\[
f_1(\lambda)u(\lambda)+f_2(\lambda)v(\lambda)=g(\lambda).
\]
又令 $f_1(\lambda)=g(\lambda)q(\lambda)$，则 $h(\lambda)=f_2(\lambda)q(\lambda)$. 对 (7.6.2) 式作初等变换，可将前两个对角元变为 $g(\lambda)$ 与 $h(\lambda)$，其余对角元不变.

现来考察 $g(\lambda)$ 与 $h(\lambda)$ 的准素因子. 将 $f_1(\lambda),f_2(\lambda)$ 作标准因式分解，其分解式不妨设为
\[
f_1(\lambda)=(\lambda-\lambda_1)^{c_1}(\lambda-\lambda_2)^{c_2}\cdots(\lambda-\lambda_t)^{c_t},
\]
\[
f_2(\lambda)=(\lambda-\lambda_1)^{d_1}(\lambda-\lambda_2)^{d_2}\cdots(\lambda-\lambda_t)^{d_t},
\]
其中 $c_i,d_i$ 为非负整数. 令
\[
e_i=\max\{c_i,d_i\},\quad k_i=\min\{c_i,d_i\},
\]
则
\[
g(\lambda)=(\lambda-\lambda_1)^{k_1}(\lambda-\lambda_2)^{k_2}\cdots(\lambda-\lambda_t)^{k_t},
\]
\[
h(\lambda)=(\lambda-\lambda_1)^{e_1}(\lambda-\lambda_2)^{e_2}\cdots(\lambda-\lambda_t)^{e_t}.
\]
不难看出 $g(\lambda),h(\lambda)$ 的准素因子组与 $f_1(\lambda),f_2(\lambda)$ 的准素因子组相同.

第二步证明 (7.6.2) 式所示矩阵的法式可通过上述变换得到.

先将第 $(1,1)$ 位置的元素依次和第 $(2,2)$ 位置，$\cdots$，第 $(n,n)$ 位置的元素进行上述变换，此时第 $(1,1)$ 元素的所有一次因式的幂都是最小的；再将第 $(2,2)$ 位置的元素依次和第 $(3,3)$ 位置，$\cdots$，第 $(n,n)$ 位置的元素进行上述变换；$\cdots$；最后将第 $(n-1,n-1)$ 位置的元素和第 $(n,n)$ 位置的元素进行上述变换. 可以看出，最后得到的对角阵就是 (7.6.2) 式所示矩阵的法式. 注意到在每一次变换的过程中，准素因子组都保持不变，这就证明了结论. $\square$

\noindent\textbf{注}\quad 引理 7.6.2 给出了求矩阵初等因子组的另外一个方法，它可以不必先求不变因子组而直接用初等变换把特征矩阵化为对角阵，再分解主对角线上的多项式即可. 另外，引理 7.6.2 的结论及其证明在一般的数域 $\mathbb{K}$ 上也成立.

\noindent\textbf{例 7.6.1}\quad 设 $\lambda I-A$ 经过初等变换后化为下列对角阵：
\[
\left(\begin{array}{ccccc}
1 & & & & \\
& (\lambda-1)^2(\lambda+2) & & & \\
& & \lambda+2 & & \\
& & & 1 & \\
& & & & \lambda-1
\end{array}\right),
\]
求 $A$ 的初等因子组.

\noindent\textbf{解}\quad 由引理 7.6.2 知，$A$ 的初等因子组为 $\lambda-1,(\lambda-1)^2,\lambda+2,\lambda+2$.

\noindent\textbf{引理 7.6.3}\quad 设 $J$ 是分块对角阵：
\[
J=\left(\begin{array}{cccc}
J_1 & & & \\
& J_2 & & \\
& & \ddots & \\
& & & J_k
\end{array}\right),
\]
其中每个 $J_i$ 都是形如引理 7.6.1 中的矩阵，$J_i$ 的初等因子组为 $(\lambda-\lambda_i)^{r_i}$，则 $J$ 的初等因子组为
\begin{equation}
(\lambda-\lambda_1)^{r_1},(\lambda-\lambda_2)^{r_2},\cdots,(\lambda-\lambda_k)^{r_k}.
\tag{7.6.3}
\end{equation}

\noindent\textbf{证明}\quad $\lambda I-J$ 是一个分块对角 $\lambda$-矩阵. 由于对分块对角阵中某一块施行初等变换时其余各块保持不变，故由引理 7.6.1 知，$\lambda I-J$ 相抵于下列分块对角阵：
\[
H=\left(\begin{array}{cccc}
H_1 & & & \\
& H_2 & & \\
& & \ddots & \\
& & & H_k
\end{array}\right),
\]
其中 $H_i=\operatorname{diag}\{1,\cdots,1,(\lambda-\lambda_i)^{r_i}\}$. 再由引理 7.6.2 即得结论. $\square$

\noindent\textbf{定理 7.6.1}\quad 设 $A$ 是复数域上的矩阵且 $A$ 的初等因子组为
\[
(\lambda-\lambda_1)^{r_1},(\lambda-\lambda_2)^{r_2},\cdots,(\lambda-\lambda_k)^{r_k},
\]
则 $A$ 相似于分块对角阵：
\begin{equation}
J=\left(\begin{array}{cccc}
J_1 & & & \\
& J_2 & & \\
& & \ddots & \\
& & & J_k
\end{array}\right),
\tag{7.6.4}
\end{equation}
其中 $J_i$ 为 $r_i$ 阶矩阵，且
\begin{equation}
J_i=\left(\begin{array}{ccccc}
\lambda_i & 1 & & & \\
& \lambda_i & 1 & & \\
& & \ddots & \ddots & \\
& & & \ddots & 1\\
& & & & \lambda_i
\end{array}\right).
\tag{7.6.5}
\end{equation}

\noindent\textbf{证明}\quad 由引理 7.6.3 知，$A$ 与 $J$ 有相同的初等因子组，因此 $A$ 与 $J$ 相似. $\square$

\noindent\textbf{定义 7.6.1}\quad (7.6.4) 式中的矩阵 $J$ 称为 $A$ 的 Jordan 标准型，每个 $J_i$ 称为 $A$ 的一个 Jordan 块.

由引理 7.6.3 可以看出，若交换任意两个 Jordan 块的位置，得到的矩阵与原来的矩阵仍有相同的初等因子组，它们仍相似. 因此矩阵 $A$ 的 Jordan 标准型中 Jordan 块的排列可以是任意的. 但是，由于每个初等因子唯一确定了一个 Jordan 块，故若不计 Jordan 块的排列次序，则矩阵的 Jordan 标准型是唯一确定的.

至此，我们对复数域上线性空间的线性变换解决了在第四章中提出的问题：求 $V$ 的一组基，使该线性变换在这组基下的表示矩阵具有简单的形式. 我们把这一结果叙述为下列定理.

\noindent\textbf{定理 7.6.2}\quad 设 $\varphi$ 是复数域上线性空间 $V$ 上的线性变换，则必存在 $V$ 的一组基，使得 $\varphi$ 在这组基下的表示矩阵为 (7.6.4) 式所示的 Jordan 标准型.

\noindent\textbf{推论 7.6.1}\quad 设 $A$ 是 $n$ 阶复矩阵，则下列结论等价：

(1) $A$ 可对角化；

(2) $A$ 的极小多项式无重根；

(3) $A$ 的初等因子都是一次多项式.

\noindent\textbf{证明}\quad (1) $\Rightarrow$ (2)：由例 6.3.2 的结论即得.

(2) $\Rightarrow$ (3)：设 $A$ 的极小多项式 $m(\lambda)$ 无重根. 由于 $m(\lambda)$ 是 $A$ 的最后一个不变因子，故 $A$ 的所有不变因子都无重根，从而 $A$ 的初等因子都是一次多项式.

(3) $\Rightarrow$ (1)：设 $A$ 的初等因子组为 $\lambda-\lambda_1,\lambda-\lambda_2,\cdots,\lambda-\lambda_n$，则由定理 7.6.1 知，$A$ 相似于对角阵 $\operatorname{diag}\{\lambda_1,\lambda_2,\cdots,\lambda_n\}$，即 $A$ 可对角化. $\square$

我们将推论 7.6.1 的几何版本叙述如下.

\noindent\textbf{推论 7.6.2}\quad 设 $\varphi$ 是复线性空间 $V$ 上的线性变换，则 $\varphi$ 可对角化当且仅当 $\varphi$ 的极小多项式无重根，当且仅当 $\varphi$ 的初等因子都是一次多项式.

\noindent\textbf{推论 7.6.3}\quad 设 $\varphi$ 是复线性空间 $V$ 上的线性变换，$V_0$ 是 $\varphi$ 的不变子空间. 若 $\varphi$ 可对角化，则 $\varphi$ 在 $V_0$ 上的限制也可对角化.

\noindent\textbf{证明}\quad 设 $\varphi,\varphi|_{V_0}$ 的极小多项式分别为 $g(\lambda),h(\lambda)$，则由推论 7.6.2 知，$g(\lambda)$ 无重根. 又 $g(\varphi|_{V_0})=g(\varphi)|_{V_0}=0$，故 $h(\lambda)\mid g(\lambda)$，于是 $h(\lambda)$ 也无重根，再次由推论 7.6.2 知，$\varphi|_{V_0}$ 可对角化. $\square$

\noindent\textbf{推论 7.6.4}\quad 设 $\varphi$ 是复线性空间 $V$ 上的线性变换，且 $V=V_1\oplus V_2\oplus\cdots\oplus V_k$，其中每个 $V_i$ 都是 $\varphi$ 的不变子空间，则 $\varphi$ 可对角化的充分必要条件是 $\varphi$ 在每个 $V_i$ 上的限制都可对角化.

\noindent\textbf{证明}\quad 必要性由推论 7.6.3 即得，下证充分性. 若 $\varphi$ 在每个 $V_i$ 上的限制都可对角化，则由定义存在 $V_i$ 的一组基，使得 $\varphi|_{V_i}$ 在这组基下的表示矩阵是对角阵. 再由定理 3.9.3 知 $V_i$ 的一组基可以拼成 $V$ 的一组基，因此 $\varphi$ 在这组基下的表示矩阵是对角阵，即 $\varphi$ 可对角化. $\square$

\noindent\textbf{推论 7.6.5}\quad 设 $A$ 是数域 $\mathbb{K}$ 上的矩阵，如果 $A$ 的特征值全在 $\mathbb{K}$ 中，则 $A$ 在 $\mathbb{K}$ 上相似于其 Jordan 标准型.

\noindent\textbf{证明}\quad 由于 $A$ 的特征值全在 $\mathbb{K}$ 中，故 $A$ 的 Jordan 标准型 $J$ 实际上是 $\mathbb{K}$ 上的矩阵. 因为 $A$ 在复数域上相似于 $J$，由推论 7.3.4 知，$A$ 在 $\mathbb{K}$ 上也相似于 $J$. $\square$

\noindent\textbf{例 7.6.2}\quad 设 $A$ 是 7 阶矩阵，其初等因子组为
\[
\lambda-1,(\lambda-1)^3,(\lambda+1)^2,\lambda-2,
\]
求其 Jordan 标准型.

\noindent\textbf{解}\quad $A$ 的 Jordan 标准型为
\[
J=\left(\begin{array}{ccccccc}
1 & & & & & & \\
& 1 & 1 & 0 & & & \\
& 0 & 1 & 1 & & & \\
& 0 & 0 & 1 & & & \\
& & & & -1 & 1 & \\
& & & & 0 & -1 & \\
& & & & & & 2
\end{array}\right),
\]
$J$ 含有 4 个 Jordan 块.

\noindent\textbf{例 7.6.3}\quad 设复数域上的四维线性空间 $V$ 上的线性变换 $\varphi$ 在一组基 $\{e_1,e_2,e_3,e_4\}$ 下的表示矩阵为
\[
A=\left(\begin{array}{cccc}
3 & 1 & 0 & 0\\
-4 & -1 & 0 & 0\\
6 & 1 & 2 & 1\\
-14 & -5 & -1 & 0
\end{array}\right),
\]
求 $V$ 的一组基，使 $\varphi$ 在这组基下的表示矩阵为 Jordan 标准型，并求出从原来的基到新基的过渡矩阵.

\noindent\textbf{解}\quad 用初等变换把 $\lambda I-A$ 化为对角 $\lambda$-矩阵并求出它的初等因子组为
\[
(\lambda-1)^2,\quad(\lambda-1)^2.
\]
因此，$A$ 的 Jordan 标准型为
\[
J=\left(\begin{array}{cccc}
1 & 1 & & \\
0 & 1 & & \\
& & 1 & 1\\
& & 0 & 1
\end{array}\right).
\]
设矩阵 $P$ 是从 $\{e_1,e_2,e_3,e_4\}$ 到新基的过渡矩阵，则
\[
P^{-1}AP=J,
\]
此即
\begin{equation}
AP=PJ.
\tag{7.6.6}
\end{equation}
设 $P=(\alpha_1,\alpha_2,\alpha_3,\alpha_4)$，其中 $\alpha_i$ 是四维列向量，代入 (7.6.6) 式得
\[
(A\alpha_1,A\alpha_2,A\alpha_3,A\alpha_4)
=(\alpha_1,\alpha_2,\alpha_3,\alpha_4)
\left(\begin{array}{cccc}
1 & 1 & & \\
0 & 1 & & \\
& & 1 & 1\\
& & 0 & 1
\end{array}\right),
\]
化成方程组为
\[
(A-I)\alpha_1=0,\quad
(A-I)\alpha_2=\alpha_1,\quad
(A-I)\alpha_3=0,\quad
(A-I)\alpha_4=\alpha_3.
\]
由于 $\alpha_1,\alpha_3$ 都是 $A$ 的属于特征值 1 的特征向量，故 $\alpha_2,\alpha_4$ 称为属于特征值 1 的广义特征向量. 我们可取方程组 $(A-I)x=0$ 的两个线性无关的解分别作为 $\alpha_1,\alpha_3$（注意不能取线性相关的两个解，因为 $P$ 是非异阵），然后再分别求出 $\alpha_2,\alpha_4$（注意诸 $\alpha_i$ 的解可能不唯一，只需取比较简单的一组解）即可. 经计算可得
\[
\alpha_1=\left(\begin{array}{c}
1\\
-2\\
1\\
-5
\end{array}\right),\quad
\alpha_2=\left(\begin{array}{c}
0\\
1\\
0\\
0
\end{array}\right),\quad
\alpha_3=\left(\begin{array}{c}
0\\
0\\
1\\
-1
\end{array}\right),\quad
\alpha_4=\left(\begin{array}{c}
0\\
0\\
0\\
1
\end{array}\right).
\]
于是
\[
P=\left(\begin{array}{cccc}
1 & 0 & 0 & 0\\
-2 & 1 & 0 & 0\\
1 & 0 & 1 & 0\\
-5 & 0 & -1 & 1
\end{array}\right).
\]
因此新基为
\[
\{e_1-2e_2+e_3-5e_4,e_2,e_3-e_4,e_4\}.
\]

\subsection*{习题 7.6}

1. 已知矩阵的下列初等因子组，写出 Jordan 标准型：

(1) $\lambda,\lambda^2,(\lambda-1)^2,(\lambda-1)^2$；

(2) $(\lambda+1)^2,\lambda-2,(\lambda-2)^3$；

(3) $(\lambda-\sqrt{2})^2,\lambda-1,(\lambda-1)^3,(\lambda-1)^3$.

2. 求非异阵 $P$，使 $P^{-1}AP$ 为 Jordan 标准型：
\[
(1)\ A=\left(\begin{array}{ccc}
0 & 1 & 0\\
-4 & 4 & 0\\
-2 & 1 & 2
\end{array}\right);\quad
(2)\ A=\left(\begin{array}{ccc}
4 & 6 & -15\\
1 & 3 & -5\\
1 & 2 & -4
\end{array}\right);
\]
\[
(3)\ A=\left(\begin{array}{ccc}
2 & -5 & 4\\
1 & -4 & 4\\
1 & -5 & 5
\end{array}\right);\quad
(4)\ A=\left(\begin{array}{cccc}
1 & -1 & 0 & 1\\
1 & 1 & 1 & 0\\
0 & -1 & 1 & 1\\
1 & 0 & 1 & 1
\end{array}\right).
\]

3. 设 $A=\operatorname{diag}\{A_1,A_2,\cdots,A_k\}$ 为分块对角阵，求证：$A$ 的初等因子组等于 $A_i$ ($i=1,\cdots,k$) 的初等因子组的无交并集. 又若交换各块的位置，则所得的矩阵仍和 $A$ 相似.

4. 设 $n$ 阶矩阵 $A$ 适合 $A^2=O$ 且 $A$ 的秩等于 $r$，试求 $A$ 的 Jordan 标准型.

5. 设 $n$ 阶矩阵 $A$ 适合 $A^2=A$ 且 $A$ 的秩等于 $r$，试求 $A$ 的 Jordan 标准型.

6. 设 $n(n>1)$ 阶矩阵 $A$ 的秩为 1，试求 $A$ 的 Jordan 标准型.

7. 设 $n(n>1)$ 阶矩阵 $A$ 的秩为 1，求证：$A$ 是幂等阵的充分必要条件是 $\operatorname{tr}(A)=1$，$A$ 是幂零阵的充分必要条件是 $\operatorname{tr}(A)=0$.

8. 设 $n$ 阶矩阵 $A$ 的极小多项式的次数等于 $n$，求证：$A$ 的 Jordan 标准型中各个 Jordan 块的主对角元素彼此不同.

9. 设 $A$ 是 $n$ 阶复矩阵且存在正整数 $k$ 使得 $A^k=I_n$，求证：$A$ 相似于对角阵.

10. 设有理数域上 $n$ 阶矩阵 $A$ 的特征多项式的所有不可约因式是 $\lambda^2+\lambda+1,\lambda^2-2$，又 $A$ 的极小多项式是四次多项式，求证：$A$ 在复数域上必相似于对角阵.

11. 设 $n$ 阶矩阵 $A$ 的全体不同特征值为 $\lambda_1,\lambda_2,\cdots,\lambda_k$，令 $g(\lambda)=(\lambda-\lambda_1)(\lambda-\lambda_2)\cdots(\lambda-\lambda_k)$，求证：$A$ 可对角化的充分必要条件是 $g(A)=O$.

12. 设 $A$ 是 $n$ 阶复矩阵，求证：$A$ 相似于分块对角阵 $\operatorname{diag}\{B,C\}$，其中 $B$ 是幂零阵，$C$ 是非异阵.

\section*{§7.7\quad Jordan 标准型的进一步讨论和应用}

在这一节里，我们要通过复线性空间 $V$ 上线性变换 $\varphi$ 的 Jordan 标准型来读出 $\varphi$ 的特征值的度数和重数，并得到全空间 $V$ 的两种直和分解，最后将给出 Jordan 标准型在矩阵理论中的一些应用.

设 $V$ 是 $n$ 维复线性空间，$\varphi$ 是 $V$ 上的线性变换. 设 $\varphi$ 的初等因子组为
\begin{equation}
(\lambda-\lambda_1)^{r_1},(\lambda-\lambda_2)^{r_2},\cdots,(\lambda-\lambda_k)^{r_k},
\tag{7.7.1}
\end{equation}
定理 7.6.2 告诉我们，存在 $V$ 的一组基 $\{e_{11},e_{12},\cdots,e_{1r_1};e_{21},e_{22},\cdots,e_{2r_2};\cdots;e_{k1},e_{k2},\cdots,e_{kr_k}\}$，使得 $\varphi$ 在这组基下的表示矩阵为
\[
J=\left(\begin{array}{cccc}
J_1 & & & \\
& J_2 & & \\
& & \ddots & \\
& & & J_k
\end{array}\right).
\]
上式中每个 $J_i$ 是相应于初等因子 $(\lambda-\lambda_i)^{r_i}$ 的 Jordan 块，其阶正好为 $r_i$. 令 $V_i$ 是由基向量 $e_{i1},e_{i2},\cdots,e_{ir_i}$ 生成的子空间，则
\begin{equation}
\begin{aligned}
\varphi(e_{i1})&=\lambda_ie_{i1},\\
\varphi(e_{i2})&=e_{i1}+\lambda_ie_{i2},\\
&\cdots\cdots\cdots\\
\varphi(e_{ir_i})&=e_{i,r_i-1}+\lambda_ie_{ir_i}.
\end{aligned}
\tag{7.7.2}
\end{equation}
这表明 $\varphi(V_i)\subseteq V_i$，即 $V_i$ ($i=1,2,\cdots,k$) 是 $\varphi$ 的不变子空间. 显然我们有
\[
V=V_1\oplus V_2\oplus\cdots\oplus V_k.
\]

线性变换 $\varphi$ 限制在 $V_1$ 上（仍记为 $\varphi$）便成为 $V_1$ 上的线性变换. 这个线性变换在基 $\{e_{11},e_{12},\cdots,e_{1r_1}\}$ 下的表示矩阵为
\[
J_1=\left(\begin{array}{ccccc}
\lambda_1 & 1 & & & \\
& \lambda_1 & 1 & & \\
& & \ddots & \ddots & \\
& & & \ddots & 1\\
& & & & \lambda_1
\end{array}\right).
\]
注意到 $J_1$ 的特征值全为 $\lambda_1$，并且 $\lambda_1I-J_1$ 的秩等于 $r_1-1$，故 $J_1$ 只有一个线性无关的特征向量，不妨选为 $e_{11}$. 显然 $e_{11}$ 也是 $\varphi$ 作为 $V$ 上线性变换关于特征值 $\lambda_1$ 的特征向量. 不失一般性，不妨设在 $\varphi$ 的初等因子组即 (7.7.1) 式中
\[
\lambda_1=\lambda_2=\cdots=\lambda_s,\quad \lambda_i\ne\lambda_1\ (i=s+1,\cdots,k);
\]
则 $J_1,\cdots,J_s$ 都以 $\lambda_1$ 为特征值，且相应于每一块有且只有一个线性无关的特征向量. 相应的特征向量可取为
\begin{equation}
e_{11},e_{21},\cdots,e_{s1},
\tag{7.7.3}
\end{equation}
显然这是 $s$ 个线性无关的特征向量. 如果 $\lambda_i\ne\lambda_1$，则容易看出 $\operatorname{r}(\lambda_1I-J_i)=r_i$，于是
\[
\operatorname{r}(\lambda_1I-J)=\sum_{i=1}^k\operatorname{r}(\lambda_1I-J_i)=(r_1-1)+\cdots+(r_s-1)+r_{s+1}+\cdots+r_k=n-s.
\]
因此 $\varphi$ 关于特征值 $\lambda_1$ 的特征子空间 $V_{\lambda_1}$ 的维数等于 $n-\operatorname{r}(\lambda_1I-J)=s$，从而特征子空间 $V_{\lambda_1}$ 以 (7.7.3) 式中的向量为一组基. 又 $\lambda_1$ 是 $\varphi$ 的 $r_1+r_2+\cdots+r_s$ 重特征值，因此 $\lambda_1$ 的重数与度数之差等于
\[
(r_1+r_2+\cdots+r_s)-s.
\]
我们把上述结论写成如下定理.

\noindent\textbf{定理 7.7.1}\quad 线性变换 $\varphi$ 的特征值 $\lambda_1$ 的度数等于 $\varphi$ 的 Jordan 标准型中属于特征值 $\lambda_1$ 的 Jordan 块的个数，$\lambda_1$ 的重数等于所有属于特征值 $\lambda_1$ 的 Jordan 块的阶数之和.

现在再来看 $J_1$ 所对应的子空间 $V_1$，由 (7.7.2) 式中诸等式可知
\[
(\varphi-\lambda_1I)(e_{1r_1})=e_{1,r_1-1},\ \cdots,\ (\varphi-\lambda_1I)(e_{12})=e_{11},\ (\varphi-\lambda_1I)(e_{11})=0,
\]
因此，若记 $\alpha=e_{1r_1}$，$\psi=\varphi-\lambda_1I$，则
\[
\psi(\alpha)=e_{1,r_1-1},\ \psi^2(\alpha)=e_{1,r_1-2},\ \cdots,\ \psi^{r_1-1}(\alpha)=e_{11},\ \psi^{r_1}(\alpha)=0.
\]
也就是说
\[
\{\alpha,\psi(\alpha),\psi^2(\alpha),\cdots,\psi^{r_1-1}(\alpha)\}
\]
构成了 $V_1$ 的一组基.

\noindent\textbf{定义 7.7.1}\quad 设 $V_0$ 是线性空间 $V$ 的 $r$ 维子空间，$\psi$ 是 $V$ 上的线性变换. 若存在 $\alpha\in V_0$，使 $\{\alpha,\psi(\alpha),\cdots,\psi^{r-1}(\alpha)\}$ 构成 $V_0$ 的一组基，则称 $V_0$ 为关于线性变换 $\psi$ 的循环子空间.

上面的事实说明，每个 Jordan 块 $J_i$ 对应的子空间 $V_i$ 是一个循环子空间. 把属于同一个特征值，比如属于 $\lambda_1$ 的所有循环子空间加起来构成 $V$ 的一个子空间：
\[
R(\lambda_1)=V_1\oplus\cdots\oplus V_s.
\]
若 $v\in R(\lambda_1)$，则不难算出 $(\varphi-\lambda_1I)^s(v)=0$，其中
\[
s=\dim R(\lambda_1)=r_1+\cdots+r_s.
\]
事实上，我们可以证明
\begin{equation}
R(\lambda_1)=\{v\in V\mid(\varphi-\lambda_1I)^n(v)=0\}.
\tag{7.7.4}
\end{equation}
为证明 (7.7.4) 式成立，设 $U=\{v\in V\mid(\varphi-\lambda_1I)^n(v)=0\}$，则由上面的分析知道，$R(\lambda_1)\subseteq U$. 另一方面，任取 $v\in U$，设 $v=v_1+v_2$，其中 $v_1\in R(\lambda_1)$，$v_2\in V_{s+1}\oplus\cdots\oplus V_k$. 因为 $(\lambda-\lambda_1)^n$ 与 $(\lambda-\lambda_{s+1})^n\cdots(\lambda-\lambda_k)^n$ 互素，故存在多项式 $p(\lambda),q(\lambda)$，使
\[
(\lambda-\lambda_1)^np(\lambda)+(\lambda-\lambda_{s+1})^n\cdots(\lambda-\lambda_k)^nq(\lambda)=1.
\]
将 $\lambda=\varphi$ 代入上式并作用在 $v$ 上可得
\[
\begin{aligned}
v&=p(\varphi)(\varphi-\lambda_1I)^n(v)+q(\varphi)(\varphi-\lambda_{s+1}I)^n\cdots(\varphi-\lambda_kI)^n(v)\\
&=q(\varphi)(\varphi-\lambda_{s+1}I)^n\cdots(\varphi-\lambda_kI)^n(v_1)+q(\varphi)(\varphi-\lambda_{s+1}I)^n\cdots(\varphi-\lambda_kI)^n(v_2)\\
&=q(\varphi)(\varphi-\lambda_{s+1}I)^n\cdots(\varphi-\lambda_kI)^n(v_1)\in R(\lambda_1).
\end{aligned}
\]
这就证明了 (7.7.4) 式.

\noindent\textbf{定义 7.7.2}\quad 设 $\lambda_0$ 是 $n$ 维复线性空间 $V$ 上线性变换 $\varphi$ 的特征值，则
\[
R(\lambda_0)=\{v\in V\mid(\varphi-\lambda_0I)^n(v)=0\}
\]
构成了 $V$ 的一个子空间，称为属于特征值 $\lambda_0$ 的根子空间.

上面的结果表明：特征值 $\lambda_0$ 的根子空间可表示为若干个循环子空间的直和，每个循环子空间对应于一个 Jordan 块.

虽然我们前面的讨论是对特征值 $\lambda_1$ 进行的，其实对任一特征值 $\lambda_i$ 均适用. 因此便有如下的定理.

\noindent\textbf{定理 7.7.2}\quad 设 $\varphi$ 是 $n$ 维复线性空间 $V$ 上的线性变换.

(1) 若 $\varphi$ 的初等因子组为
\[
(\lambda-\lambda_1)^{r_1},(\lambda-\lambda_2)^{r_2},\cdots,(\lambda-\lambda_k)^{r_k},
\]
则 $V$ 可分解为 $k$ 个不变子空间的直和：
\begin{equation}
V=V_1\oplus V_2\oplus\cdots\oplus V_k,
\tag{7.7.5}
\end{equation}
其中 $V_i$ 是维数等于 $r_i$ 的关于 $\varphi-\lambda_iI$ 的循环子空间；

(2) 若 $\lambda_1,\cdots,\lambda_s$ 是 $\varphi$ 的全体不同特征值，则 $V$ 可分解为 $s$ 个不变子空间的直和：
\begin{equation}
V=R(\lambda_1)\oplus R(\lambda_2)\oplus\cdots\oplus R(\lambda_s),
\tag{7.7.6}
\end{equation}
其中 $R(\lambda_i)$ 是 $\lambda_i$ 的根子空间，$R(\lambda_i)$ 的维数等于 $\lambda_i$ 的重数，且每个 $R(\lambda_i)$ 又可分解为 (7.7.5) 式中若干个 $V_j$ 的直和.

下面我们举例说明 Jordan 标准型在矩阵理论中的应用.

\noindent\textbf{例 7.7.1}\quad 证明：复数域上的方阵 $A$ 必可分解为两个对称阵的乘积.

\noindent\textbf{证明}\quad 设 $P$ 是非异阵且使 $P^{-1}AP=J$ 为 $A$ 的 Jordan 标准型，于是 $A=PJP^{-1}$. 设 $J_i$ 是 $J$ 的第 $i$ 个 Jordan 块，则
\[
J_i=\left(\begin{array}{ccccc}
\lambda_i & 1 & & & \\
& \lambda_i & 1 & & \\
& & \ddots & \ddots & \\
& & & \ddots & 1\\
& & & & \lambda_i
\end{array}\right)
=\left(\begin{array}{ccccc}
& & & & \lambda_i\\
& & & \lambda_i & 1\\
& & \ddots & \ddots & \\
& \lambda_i & 1 & & \\
\lambda_i & 1 & & &
\end{array}\right)
\left(\begin{array}{ccccc}
& & & & 1\\
& & & 1 & \\
& & \ddots & & \\
& 1 & & & \\
1 & & & &
\end{array}\right),
\]
即 $J_i$ 可分解为两个对称阵之积. 因此 $J$ 也可以分解为两个对称阵之积，记为 $S_1S_2$，于是
\[
A=PJP^{-1}=PS_1S_2P^{-1}=(PS_1P')((P^{-1})'S_2P^{-1}).\quad \square
\]

\noindent\textbf{例 7.7.2}\quad 设 $A$ 是本章例 7.6.3 中的矩阵，计算 $A^k$ ($k\geq 1$).

\noindent\textbf{解}\quad 因为
\[
P^{-1}A^kP=(P^{-1}AP)^k=J^k,
\]
故先计算 $J^k$. 注意 $J$ 是分块对角阵，它的 $k$ 次方等于将各对角块 $k$ 次方，因此
\[
J^k=\left(\begin{array}{cccc}
1 & 1 & & \\
0 & 1 & & \\
& & 1 & 1\\
& & 0 & 1
\end{array}\right)^k
=\left(\begin{array}{cccc}
1 & k & & \\
0 & 1 & & \\
& & 1 & k\\
& & 0 & 1
\end{array}\right),
\]
\[
\begin{aligned}
A^k=PJ^kP^{-1}
&=\left(\begin{array}{cccc}
1 & 0 & 0 & 0\\
-2 & 1 & 0 & 0\\
1 & 0 & 1 & 0\\
-5 & 0 & -1 & 1
\end{array}\right)
\left(\begin{array}{cccc}
1 & k & & \\
0 & 1 & & \\
& & 1 & k\\
& & 0 & 1
\end{array}\right)
\left(\begin{array}{cccc}
1 & 0 & 0 & 0\\
-2 & 1 & 0 & 0\\
1 & 0 & 1 & 0\\
-5 & 0 & -1 & 1
\end{array}\right)^{-1}\\
&=\left(\begin{array}{cccc}
2k+1 & k & 0 & 0\\
-4k & -2k+1 & 0 & 0\\
6k & k & k+1 & k\\
-14k & -5k & -k & -k+1
\end{array}\right).
\end{aligned}
\]

下面我们要用 Jordan 标准型来证明著名的 Jordan-Chevalley（若当-谢瓦莱）分解定理，它在 Lie（李）代数中有重要的应用. 为此我们先证明一个引理.

\noindent\textbf{引理 7.7.1}\quad 设 $A,B$ 是两个 $n$ 阶可对角化复矩阵且 $AB=BA$，则它们可同时对角化，即存在可逆阵 $P$，使 $P^{-1}AP$ 和 $P^{-1}BP$ 都是对角阵.

\noindent\textbf{证明}\quad 转化为几何的语言：设 $\varphi,\psi$ 是 $n$ 维复线性空间 $V$ 上两个可对角化线性变换且 $\varphi\psi=\psi\varphi$，则它们可同时对角化，即存在 $V$ 的一组基，使 $\varphi,\psi$ 在这组基下的表示矩阵都是对角阵. 设 $\lambda_1,\lambda_2,\cdots,\lambda_s$ 是 $\varphi$ 的全体不同特征值，$V_i$ ($i=1,2,\cdots,s$) 是属于特征值 $\lambda_i$ 的特征子空间. 若 $s=1$，则由 $\varphi$ 可对角化不难推出 $\varphi=\lambda_1I_V$ 是数量变换. 此时可取 $V$ 的一组基，使 $\psi$ 的表示矩阵是对角阵，则 $\varphi$ 的表示矩阵是 $I_n$，结论显然成立. 以下不妨设 $s>1$，由 $\varphi$ 可对角化知
\[
V=V_1\oplus V_2\oplus\cdots\oplus V_s.
\]
由于 $\varphi\psi=\psi\varphi$，故对任意的 $\alpha\in V_i$，
\[
\varphi\psi(\alpha)=\psi\varphi(\alpha)=\lambda_i\psi(\alpha),
\]
这表明 $\psi(\alpha)\in V_i$，即 $V_i$ 是 $\psi$ 的不变子空间. 将 $\varphi$ 和 $\psi$ 限制在 $V_i$ 上，此时 $\varphi|_{V_i}=\lambda_iI_{V_i}$ 是数量变换，由推论 7.6.3 知 $\psi|_{V_i}$ 仍是可对角化线性变换. 由 $s=1$ 的情形，存在 $V_i$ 的一组基，使 $\varphi|_{V_i},\psi|_{V_i}$ 在这组基下的表示矩阵都是对角阵. 将各个 $V_i$ 的基合并成 $V$ 的一组基，则 $\varphi,\psi$ 在这组基下的表示矩阵都是对角阵. $\square$

\noindent\textbf{定理 7.7.3 (Jordan-Chevalley 分解)}\quad 设 $A$ 是 $n$ 阶复矩阵，则 $A$ 可分解为 $A=B+C$，其中 $B,C$ 适合下面条件：

(1) $B$ 是一个可对角化矩阵；

(2) $C$ 是一个幂零阵；

(3) $BC=CB$；

(4) $B,C$ 均可表示为 $A$ 的多项式.

不仅如此，上述满足条件 (1)～(3) 的分解是唯一的.

\noindent\textbf{证明}\quad 先对 $A$ 的 Jordan 标准型 $J$ 证明结论. 设 $A$ 的全体不同特征值为 $\lambda_1,\lambda_2,\cdots,\lambda_s$ 且
\[
J=\left(\begin{array}{cccc}
J_1 & & & \\
& J_2 & & \\
& & \ddots & \\
& & & J_s
\end{array}\right),
\]
其中 $J_i$ 是属于特征值 $\lambda_i$ 的根子空间对应的块，其阶设为 $m_i$. 显然对每个 $i$ 均有 $J_i=M_i+N_i$，其中 $M_i=\lambda_iI$ 是对角阵，$N_i$ 是幂零阵且 $M_iN_i=N_iM_i$. 令
\[
M=\left(\begin{array}{cccc}
M_1 & & & \\
& M_2 & & \\
& & \ddots & \\
& & & M_s
\end{array}\right),\quad
N=\left(\begin{array}{cccc}
N_1 & & & \\
& N_2 & & \\
& & \ddots & \\
& & & N_s
\end{array}\right),
\]
则 $J=M+N$，$MN=NM$，$M$ 是对角阵，$N$ 是幂零阵.

因为 $(J_i-\lambda_iI)^{m_i}=0$，所以 $J_i$ 适合多项式 $(\lambda-\lambda_i)^{m_i}$. 而 $\lambda_i$ 互不相同，因此多项式 $(\lambda-\lambda_1)^{m_1},(\lambda-\lambda_2)^{m_2},\cdots,(\lambda-\lambda_s)^{m_s}$ 两两互素. 由中国剩余定理，存在多项式 $g(\lambda)$ 满足条件
\[
g(\lambda)=h_i(\lambda)(\lambda-\lambda_i)^{m_i}+\lambda_i,
\]
对所有 $i=1,2,\cdots,s$ 成立（这里 $h_i(\lambda)$ 也是多项式）. 代入 $J_i$ 得到
\[
g(J_i)=h_i(J_i)(J_i-\lambda_iI)^{m_i}+\lambda_iI=\lambda_iI=M_i.
\]
于是
\[
g(J)=\left(\begin{array}{cccc}
g(J_1) & & & \\
& g(J_2) & & \\
& & \ddots & \\
& & & g(J_s)
\end{array}\right)
=\left(\begin{array}{cccc}
M_1 & & & \\
& M_2 & & \\
& & \ddots & \\
& & & M_s
\end{array}\right)=M.
\]
又因为 $N=J-M=J-g(J)$，所以 $N$ 也是 $J$ 的多项式.

现考虑一般情形，设 $P^{-1}AP=J$，则 $A=PJP^{-1}=P(M+N)P^{-1}$. 令 $B=PMP^{-1}$，$C=PNP^{-1}$，则 $B$ 是可对角化矩阵，$C$ 是幂零阵，$BC=CB$ 并且
\[
g(A)=g(PJP^{-1})=Pg(J)P^{-1}=PMP^{-1}=B,
\]
从而 $C=A-g(A)$.

最后证明唯一性. 假设 $A$ 有另一满足条件 (1)～(3) 的分解 $A=B_1+C_1$，则 $B-B_1=C_1-C$. 由 $B_1C_1=C_1B_1$ 不难验证 $AB_1=B_1A$，$AC_1=C_1A$. 因为 $B=g(A)$，故 $BB_1=B_1B$. 同理 $CC_1=C_1C$. 设 $C^r=O$，$C_1^t=O$，用二项式定理即知 $(C_1-C)^{r+t}=O$. 于是
\[
(B-B_1)^{r+t}=(C_1-C)^{r+t}=O.
\]
因为 $BB_1=B_1B$，它们都是可对角化矩阵，由引理知它们可同时对角化，即存在可逆阵 $Q$，使 $Q^{-1}BQ$ 和 $Q^{-1}B_1Q$ 都是对角阵. 注意到
\[
(Q^{-1}BQ-Q^{-1}B_1Q)^{r+t}=(Q^{-1}(B-B_1)Q)^{r+t}=Q^{-1}(B-B_1)^{r+t}Q=O,
\]
两个对角阵之差仍是一个对角阵，这个差的幂零等于零矩阵，则这两个矩阵必相等，由此即得 $B=B_1$，从而 $C=C_1$. $\square$

\subsection*{习题 7.7}

1. 求 $n$ 阶 Jordan 块
\[
J=\left(\begin{array}{ccccc}
\lambda_0 & 1 & & & \\
& \lambda_0 & 1 & & \\
& & \ddots & \ddots & \\
& & & \ddots & 1\\
& & & & \lambda_0
\end{array}\right)
\]
的 $k$ ($k\geq 1$) 次幂 $J^k$.

2. 求下列矩阵的 $k$ ($k\geq 1$) 次幂：
\[
(1)\ \left(\begin{array}{rrrr}
3 & -4 & 0 & 2\\
4 & -5 & -2 & 4\\
0 & 0 & 3 & -2\\
0 & 0 & 2 & -1
\end{array}\right);\quad
(2)\ \left(\begin{array}{rrrr}
4 & -1 & 1 & -7\\
9 & -2 & -7 & -1\\
0 & 0 & 5 & -8\\
0 & 0 & 2 & -3
\end{array}\right).
\]

3. 求矩阵 $B$，使 $A=B^2$，其中
\[
A=\left(\begin{array}{rr}
3 & 1\\
-1 & 5
\end{array}\right).
\]

4. 证明：例 7.7.1 中把矩阵分解为两个对称阵之积时，可指定其中任意一个矩阵为非异阵.

5. $n$ 阶矩阵 $A$ 的特征值全为 1 且只有一个线性无关的特征向量，求 $A$ 的 Jordan 标准型.

6. 求证：$n$ 阶方阵 $A$ 的秩为 $r$ 的充分必要条件是 $A$ 的形如 $\lambda^k$ 的初等因子恰有 $n-r$ 个.

7. 设 $\lambda_0$ 是 $n$ 阶矩阵 $A$ 的 $k$ 重特征值，求证：$\operatorname{r}((\lambda_0I_n-A)^k)=n-k$.

8. 设 $n$ 阶复矩阵 $A$ 有特征值 0，且满足 $\operatorname{r}(A^k)=\operatorname{r}(A^{k+1})$. 若 $\lambda^m$ 是 $A$ 的一个初等因子，求证：$m\leq k$.

9. 求下列矩阵的 Jordan 标准型：
\[
A=\left(\begin{array}{cccccc}
1 & 2 & 3 & 4 & \cdots & n\\
0 & 1 & 2 & 3 & \cdots & n-1\\
0 & 0 & 1 & 2 & \cdots & n-2\\
\vdots & \vdots & \vdots & \vdots & \ddots & \vdots\\
0 & 0 & 0 & 0 & \cdots & 1
\end{array}\right).
\]

10. 求下列矩阵的 Jordan 标准型，其中 $a$ 为参数：
\[
A=\left(\begin{array}{rrrr}
1 & a & 0 & 2\\
0 & 1 & 0 & -1\\
-3 & 4 & 1 & 3\\
0 & 0 & 0 & 1
\end{array}\right).
\]

11. 设 $n$ 阶矩阵 $A$ 的特征值全为 1，求证：对任意的正整数 $k$，$A^k$ 与 $A$ 相似.

12. 设 $n$ 阶矩阵 $A$ 的特征值全为 1 或 $-1$，求证：$A^{-1}$ 与 $A$ 相似.

\section*{* §7.8\quad 矩阵函数}

在这一节里我们将要定义矩阵函数，如 $e^A,\sin A,\cos A$ 等，方法是利用幂级数来定义这些函数.

我们曾经遇到过矩阵多项式：
\[
f(A)=a_mA^m+a_{m-1}A^{m-1}+\cdots+a_0I.
\]
让 $A$ 在 $M_n(\mathbb C)$（即 $n$ 阶复矩阵集合）上变动，$f(A)$ 就成了矩阵函数（值也在 $M_n(\mathbb C)$ 中）. 对矩阵多项式函数，可以用 Jordan 标准型来简化它的计算. 设 $P$ 是非异阵，使
\[
P^{-1}AP=J=\operatorname{diag}\{J_1,J_2,\cdots,J_k\}
\]
是 $A$ 的 Jordan 标准型，其中 $J_i$ 是 Jordan 块，则
\[
J^m=\operatorname{diag}\{J_1^m,J_2^m,\cdots,J_k^m\}.
\]
又
\[
A^m=(PJP^{-1})^m=PJ^mP^{-1},
\]
因此要计算 $f(A)$，只需计算出 $J_i^m$ 即可. 设 $J_i$ 是特征值为 $\lambda_i$ 的 $r$ 阶 Jordan 块：
\[
J_i=\left(\begin{array}{ccccc}
\lambda_i & 1 & & & \\
& \lambda_i & 1 & & \\
& & \ddots & \ddots & \\
& & & \ddots & 1\\
& & & & \lambda_i
\end{array}\right),
\]
用数学归纳法不难证明
\begin{equation}
J_i^m=\left(\begin{array}{ccccc}
\lambda_i^m & C_m^1\lambda_i^{m-1} & C_m^2\lambda_i^{m-2} & \cdots & \cdots\\
& \lambda_i^m & C_m^1\lambda_i^{m-1} & \cdots & \cdots\\
& & \lambda_i^m & \cdots & \cdots\\
& & & \ddots & \vdots\\
& & & & \lambda_i^m
\end{array}\right).
\tag{7.8.1}
\end{equation}
若
\[
f(x)=a_0+a_1x+\cdots+a_px^p,
\]
则不难算出
\begin{equation}
f(J_i)=\left(\begin{array}{ccccc}
f(\lambda_i) & \dfrac1{1!}f'(\lambda_i) & \dfrac1{2!}f^{(2)}(\lambda_i) & \cdots & \dfrac1{(r-1)!}f^{(r-1)}(\lambda_i)\\
& f(\lambda_i) & \dfrac1{1!}f'(\lambda_i) & \cdots & \dfrac1{(r-2)!}f^{(r-2)}(\lambda_i)\\
& & f(\lambda_i) & \cdots & \dfrac1{(r-3)!}f^{(r-3)}(\lambda_i)\\
& & & \ddots & \vdots\\
& & & & f(\lambda_i)
\end{array}\right).
\tag{7.8.2}
\end{equation}
再由
\[
\begin{aligned}
f(A)&=f(PJP^{-1})=Pf(J)P^{-1}\\
&=Pf(\operatorname{diag}\{J_1,J_2,\cdots,J_k\})P^{-1}\\
&=P\operatorname{diag}\{f(J_1),f(J_2),\cdots,f(J_k)\}P^{-1},
\end{aligned}
\]
即可计算出 $f(A)$.

现在引进 $n$ 阶复方阵幂级数的概念. 首先要定义何为方阵幂级数的收敛？设有 $n$ 阶复方阵序列 $\{A_p\}$：
\[
A_p=\left(\begin{array}{ccc}
a_{11}^{(p)} & \cdots & a_{1n}^{(p)}\\
\vdots & & \vdots\\
a_{n1}^{(p)} & \cdots & a_{nn}^{(p)}
\end{array}\right),
\]
$B=(b_{ij})$ 是一个同阶方阵，若对每个 $(i,j)$，序列 $\{a_{ij}^{(p)}\}$ 均收敛于 $b_{ij}$，即
\[
\lim_{p\to\infty}a_{ij}^{(p)}=b_{ij},
\]
则称矩阵序列 $\{A_p\}$ 收敛于 $B$，记为
\[
\lim_{p\to\infty}A_p=B.
\]
否则称 $\{A_p\}$ 发散. 设
\[
f(z)=a_0+a_1z+a_2z^2+\cdots
\]
是一个幂级数，记
\[
f_p(z)=a_0+a_1z+a_2z^2+\cdots+a_pz^p
\]
是其部分和. 若矩阵序列 $\{f_p(A)\}$ 收敛于 $B$，则称矩阵级数
\begin{equation}
f(A)=a_0I+a_1A+a_2A^2+\cdots
\tag{7.8.3}
\end{equation}
收敛，极限为 $B$，记为 $f(A)=B$. 否则称 $f(A)$ 发散. 用变量矩阵 $X$ 代替 $A$，便可定义矩阵幂级数
\begin{equation}
f(X)=a_0I+a_1X+a_2X^2+\cdots.
\tag{7.8.4}
\end{equation}

\noindent\textbf{定理 7.8.1}\quad 设 $f(z)=\sum_{i=0}^{\infty}a_iz^i$ 是复幂级数，则

(1) 方阵幂级数 $f(X)$ 收敛的充分必要条件是对任一非异阵 $P$，$f(P^{-1}XP)$ 都收敛，这时
\[
f(P^{-1}XP)=P^{-1}f(X)P;
\]

(2) 若 $X=\operatorname{diag}\{X_1,\cdots,X_k\}$，则 $f(X)$ 收敛的充分必要条件是 $f(X_1),\cdots,f(X_k)$ 都收敛，这时
\[
f(X)=\operatorname{diag}\{f(X_1),\cdots,f(X_k)\};
\]

(3) 若 $f(z)$ 的收敛半径为 $r$，$J_0$ 是特征值为 $\lambda_0$ 的 $n$ 阶 Jordan 块
\[
J_0=\left(\begin{array}{ccccc}
\lambda_0 & 1 & & & \\
& \lambda_0 & 1 & & \\
& & \ddots & \ddots & \\
& & & \ddots & 1\\
& & & & \lambda_0
\end{array}\right),
\]
则当 $|\lambda_0|<r$ 时 $f(J_0)$ 收敛，且
\[
f(J_0)=\left(\begin{array}{ccccc}
f(\lambda_0) & \dfrac1{1!}f'(\lambda_0) & \dfrac1{2!}f^{(2)}(\lambda_0) & \cdots & \dfrac1{(n-1)!}f^{(n-1)}(\lambda_0)\\
& f(\lambda_0) & \dfrac1{1!}f'(\lambda_0) & \cdots & \dfrac1{(n-2)!}f^{(n-2)}(\lambda_0)\\
& & f(\lambda_0) & \cdots & \dfrac1{(n-3)!}f^{(n-3)}(\lambda_0)\\
& & & \ddots & \vdots\\
& & & & f(\lambda_0)
\end{array}\right).
\]

\noindent\textbf{证明}\quad 设 $f_p(z)=a_0+a_1z+a_2z^2+\cdots+a_pz^p$ 是 $f(z)$ 前 $p+1$ 项的部分和.

(1) 注意到 $f_p(z)$ 是多项式，从而有
\[
f_p(P^{-1}XP)=P^{-1}f_p(X)P.
\]
由于 $n$ 阶矩阵序列的收敛等价于 $n^2$ 个数值序列的收敛，故
\[
\begin{aligned}
f(P^{-1}XP)&=\lim_{p\to\infty}f_p(P^{-1}XP)=\lim_{p\to\infty}P^{-1}f_p(X)P\\
&=P^{-1}\left(\lim_{p\to\infty}f_p(X)\right)P=P^{-1}f(X)P.
\end{aligned}
\]

(2) 注意到 $f_p(z)$ 是多项式，从而有
\[
f_p(X)=f_p(\operatorname{diag}\{X_1,\cdots,X_k\})=\operatorname{diag}\{f_p(X_1),\cdots,f_p(X_k)\}.
\]
由于分块矩阵序列的收敛等价于每个分块的矩阵序列的收敛，故
\[
\begin{aligned}
f(X)&=\lim_{p\to\infty}f_p(X)=\lim_{p\to\infty}\operatorname{diag}\{f_p(X_1),\cdots,f_p(X_k)\}\\
&=\operatorname{diag}\left\{\lim_{p\to\infty}f_p(X_1),\cdots,\lim_{p\to\infty}f_p(X_k)\right\}=\operatorname{diag}\{f(X_1),\cdots,f(X_k)\}.
\end{aligned}
\]

(3) 由本节开始部分的计算可知
\begin{equation}
f_p(J_0)=\left(\begin{array}{ccccc}
f_p(\lambda_0) & \dfrac1{1!}f_p'(\lambda_0) & \dfrac1{2!}f_p^{(2)}(\lambda_0) & \cdots & \dfrac1{(n-1)!}f_p^{(n-1)}(\lambda_0)\\
& f_p(\lambda_0) & \dfrac1{1!}f_p'(\lambda_0) & \cdots & \dfrac1{(n-2)!}f_p^{(n-2)}(\lambda_0)\\
& & f_p(\lambda_0) & \cdots & \dfrac1{(n-3)!}f_p^{(n-3)}(\lambda_0)\\
& & & \ddots & \vdots\\
& & & & f_p(\lambda_0)
\end{array}\right).
\tag{7.8.5}
\end{equation}
令 $p\to\infty$，由矩阵序列收敛与 $n^2$ 个数值序列收敛的等价性即得结论. $\square$

定理 7.8.1 使我们可以利用 Jordan 标准型来简化矩阵幂级数的计算以及矩阵幂级数收敛性的讨论.

\noindent\textbf{定理 7.8.2}\quad 设 $f(z)$ 是复幂级数，收敛半径为 $r$. 设 $A$ 是 $n$ 阶复方阵，特征值为 $\lambda_1,\lambda_2,\cdots,\lambda_n$，定义 $A$ 的谱半径
\[
\rho(A)=\max_{1\leq i\leq n}|\lambda_i|.
\]

(1) 若 $\rho(A)<r$，则 $f(A)$ 收敛；

(2) 若 $\rho(A)>r$，则 $f(A)$ 发散；

(3) 若 $\rho(A)=r$，则 $f(A)$ 收敛的充分必要条件是：对每一个模长等于 $r$ 的特征值 $\lambda_j$，若 $A$ 的属于 $\lambda_j$ 的初等因子中最高幂为 $n_j$ 次，则 $n_j$ 个数值级数
\begin{equation}
f(\lambda_j),\ f'(\lambda_j),\ \cdots,\ f^{(n_j-1)}(\lambda_j)
\tag{7.8.6}
\end{equation}
都收敛；

(4) 若 $f(A)$ 收敛，则 $f(A)$ 的特征值为
\[
f(\lambda_1),\ f(\lambda_2),\ \cdots,\ f(\lambda_n).
\]

\noindent\textbf{证明}\quad 设 $A$ 的 Jordan 标准型为 $J=\operatorname{diag}\{J_1,J_2,\cdots,J_k\}$. 显然 $f(A)$ 的收敛性等价于所有 $f(J_i)$ ($i=1,\cdots,k$) 的收敛性. 由上面的定理即知 (1) 成立.

若某一个 $|\lambda_j|>r$，则 $f(\lambda_j)$ 发散，因此 $f(J_j)$ 发散，故 $f(A)$ 发散，这就证明了 (2).

当 $\rho(A)=r$ 时，对 $|\lambda_i|<r$ 的 $J_i$，$f(J_i)$ 收敛. 对 $|\lambda_j|=r$ 的特征值 $\lambda_j$，注意到 $f(z)$ 的任意次导数的收敛半径仍为 $r$，又初等因子 $(\lambda-\lambda_j)^{n_j}$ 对应的 Jordan 块为 $n_j$ 阶，从 (7.8.5) 式即可知道 $f(J_j)$ 的收敛性等价于 (7.8.6) 式中 $n_j$ 个级数的收敛性.

最后若 $f(A)$ 收敛，则 $f(A)$ 与 $f(J)$ 有相同的特征值，即为 $f(\lambda_1),f(\lambda_2),\cdots,f(\lambda_n)$. $\square$

由复分析知道：
\[
\begin{aligned}
e^z&=1+\frac1{1!}z+\frac1{2!}z^2+\frac1{3!}z^3+\cdots,\\
\sin z&=z-\frac1{3!}z^3+\frac1{5!}z^5-\frac1{7!}z^7+\cdots,\\
\cos z&=1-\frac1{2!}z^2+\frac1{4!}z^4-\frac1{6!}z^6+\cdots,\\
\ln(1+z)&=z-\frac12z^2+\frac13z^3-\frac14z^4+\cdots.
\end{aligned}
\]
前 3 个级数在整个复平面上收敛，而 $\ln(1+z)$ 的收敛半径为 1，于是对一切方阵，定义
\[
\begin{aligned}
e^A&=I+\frac1{1!}A+\frac1{2!}A^2+\frac1{3!}A^3+\cdots,\\
\sin A&=A-\frac1{3!}A^3+\frac1{5!}A^5-\frac1{7!}A^7+\cdots,\\
\cos A&=I-\frac1{2!}A^2+\frac1{4!}A^4-\frac1{6!}A^6+\cdots
\end{aligned}
\]
都有意义. 若 $A$ 所有特征值的模长都小于 1，则
\[
\ln(I+A)=A-\frac12A^2+\frac13A^3-\frac14A^4+\cdots
\]
也有意义. 同理还可以定义幂函数、双曲函数等. 矩阵函数在分析数学中，比如在微分方程理论中有着重要的应用.

在具体计算矩阵函数时必须注意不要随便套用数值函数的性质，比如在数值函数中，成立
\[
e^x\cdot e^y=e^{x+y},
\]
但一般来说对矩阵 $A,B$，下面的等式并不一定成立：
\[
e^A\cdot e^B=e^{A+B}.
\]
如对
\[
A=\left(\begin{array}{cc}
1 & 1\\
0 & 0
\end{array}\right),\quad
B=\left(\begin{array}{cc}
0 & 0\\
1 & 1
\end{array}\right),
\]
读者可验证 $e^A\cdot e^B\ne e^{A+B}$. 但如果 $A$ 与 $B$ 乘法可交换，即 $AB=BA$，则 $e^A\cdot e^B=e^{A+B}$ 必成立. 我们把这一结论作为习题留给读者自己证明.

我们举例来说明如何计算 $e^{tA}$，其中 $t$ 是一个数值变量，$A$ 是一个 $n$ 阶复方阵. 若令 $f(z)=e^{tz}$，则由定理 7.8.1 即得 $f(A)=e^{tA}$ 的计算结果. 下面我们换一种方法来计算. 先设 $J_i$ 是特征值为 $\lambda_i$ 的 $r$ 阶 Jordan 块，则
\[
J_i=\lambda_iI+N,
\]
其中 $N$ 是 $r$ 阶幂零阵，即 $N^r=O$，于是
\[
e^N=I+N+\frac1{2!}N^2+\frac1{3!}N^3+\cdots+\frac1{(r-1)!}N^{r-1}.
\]
因为 $(\lambda_iI)N=N(\lambda_iI)$，故
\[
\begin{aligned}
e^{J_i}&=e^{\lambda_iI+N}=e^{\lambda_iI}\cdot e^N=e^{\lambda_i}\cdot e^N\\
&=e^{\lambda_i}I+e^{\lambda_i}N+\frac1{2!}e^{\lambda_i}N^2+\cdots+\frac1{(r-1)!}e^{\lambda_i}N^{r-1}.
\end{aligned}
\]
同理
\[
e^{tJ_i}=e^{t\lambda_i}\cdot e^{tN}=e^{t\lambda_i}\left(\begin{array}{ccccc}
1 & t & \dfrac1{2!}t^2 & \dfrac1{3!}t^3 & \cdots\ \dfrac1{(r-1)!}t^{r-1}\\
& 1 & t & \dfrac1{2!}t^2 & \cdots\ \dfrac1{(r-2)!}t^{r-2}\\
& & 1 & t & \cdots\ \dfrac1{(r-3)!}t^{r-3}\\
& & & \ddots & \vdots\\
& & & & 1
\end{array}\right).
\]
现设 $P^{-1}AP=J=\operatorname{diag}\{J_1,\cdots,J_k\}$ 是 $A$ 的 Jordan 标准型，则
\[
e^{tA}=e^{P(tJ)P^{-1}}=Pe^{tJ}P^{-1},
\]
\[
e^{tJ}=\left(\begin{array}{ccc}
e^{tJ_1} & & \\
& \ddots & \\
& & e^{tJ_k}
\end{array}\right).
\]
将 $e^{tJ_i}$ 的式子代入上面的式子即可求出 $e^{tA}$.

\subsection*{习题 7.8}

1. 若 $n$ 阶矩阵 $A$ 和 $B$ 乘法可交换，求证：
\[
e^A\cdot e^B=e^B\cdot e^A.
\]

2. 若 $n$ 阶矩阵 $A$ 和 $B$ 乘法可交换，求证：
\[
e^A\cdot e^B=e^{A+B}.
\]

3. 设 $A$ 是 $n$ 阶矩阵，求证：$\sin^2A+\cos^2A=I_n$.

4. 设 $A$ 是 $n$ 阶矩阵，求证：$\sin2A=2\sin A\cos A$.

5. 计算 $\sin(e^{cI})$ 及 $\cos(e^{cI})$，其中 $c$ 是非零常数.

6. 设 $A$ 是 $n$ 阶方阵，求 $e^A$ 的行列式.

7. 求证：对任一 $n$ 阶方阵 $A$，$e^A$ 总是非异阵.

8. 设 $A$ 是 $n$ 阶方阵，求 $\lim_{k\to\infty}A^k$ 存在的充分必要条件以及极限矩阵.

\section*{历史与展望}

历史上，对行列式和矩阵的研究并非只限定在数域（如常见的实数域 $\mathbb R$ 和复数域 $\mathbb C$）上，在带有加法和乘法两种运算的某一类环（如常见的整数环 $\mathbb Z$ 和一元多项式环 $\mathbb K[\lambda]$）上，数学家们也引入了许多新的概念，获得了很多重要的结果.

设某一类环上的矩阵 $A,B$ 相抵（等价），即存在这类环上的可逆阵 $P,Q$，使 $B=PAQ$. Sylvester 在 1851 年关于行列式的论文中证明了：$B$ 的 $i$ 阶子式的最大公因子 $D_i'$ 等于 $A$ 的 $i$ 阶子式的最大公因子 $D_i$. 这些 $D_i$ 称为 $A$ 的行列式因子. 1861 年 Smith（史密斯）在研究整数矩阵时证明了：每个秩为 $r$ 的矩阵 $A$ 相抵于一个对角阵 $\operatorname{diag}\{d_1,\cdots,d_r,0,\cdots,0\}$，其中 $d_i\mid d_{i+1}$ ($1\leq i\leq r-1$). 这些 $d_i$ 称为 $A$ 的不变因子，这一对角阵后来被称为 Smith 标准型. 由 Sylvester 的结论，即得行列式因子与不变因子之间的互推关系式：
\[
d_1=D_1,\quad d_2=D_2/D_1,\quad \cdots,\quad d_r=D_r/D_{r-1}.
\]

进一步，若设 $d_i=p_1^{l_{i1}}p_2^{l_{i2}}\cdots p_k^{l_{ik}}$，其中 $p_i$ 是素数（素元），则这些 $p_j^{l_{ij}}$ 称为 $A$ 的初等因子. 不变因子确定初等因子，反之亦然. 不变因子和初等因子的意义在于：某一类环上的矩阵 $A$ 和 $B$ 相抵当且仅当 $A$ 和 $B$ 有相同的不变因子或初等因子.

不变因子和初等因子的概念，其实最早是从 Sylvester 和 Weierstrass 的行列式的工作中产生，后来由 Frobenius 在 1878 年的论文中将它们引入到矩阵中的. Frobenius 在那篇论文中对不变因子做了进一步的工作，然后以合乎逻辑的形式整理了不变因子和初等因子的理论，他还用“逆步变换”的名称处理了矩阵 $A$ 到 $B$ 的相似变换. 后世为了纪念 Frobenius 在相似标准型理论创建过程中的重要贡献，将基于不变因子的有理标准型也称为 Frobenius 标准型.

利用初等因子的概念，法国数学家 Jordan 在 1870 年证明了：任一复矩阵均相似于分块对角阵 $J=\operatorname{diag}\{J_{r_1}(\lambda_1),J_{r_2}(\lambda_2),\cdots,J_{r_k}(\lambda_k)\}$，其中 $J_{r_i}(\lambda_i)$ 是一个上三角阵，主对角线上元素全为 $\lambda_i$，上次对角线上元素全为 1，其余元素全为 0. 这一分块对角阵称为 Jordan 标准型，主对角块 $J_{r_i}(\lambda_i)$ 称为 Jordan 块.

Metzler（梅茨勒）在 1892 年的论文中引入了矩阵函数的概念. 他把矩阵的超越函数写成矩阵的幂级数，并对 $e^M,e^{-M},\ln M,\sin M$ 和 $\arcsin M$ 建立了级数，例如 $e^M=\sum_{n=0}^{\infty}M^n/n!$.

\section*{复习题七}

1. 设 $n$ 阶方阵 $A,B,C,D$ 中 $A,C$ 可逆，求证：存在可逆阵 $P,Q$，使得 $A=PCQ$，$B=PDQ$ 的充分必要条件是 $\lambda A-B$ 与 $\lambda C-D$ 相抵.

2. 求证：存在 $n$ 阶实方阵 $A$，满足 $A^2+2A+5I_n=O$ 的充分必要条件是 $n$ 为偶数. 当 $n\geq 4$ 时，验证满足上述条件的矩阵 $A$ 有无限个不变子空间.

3. 设 $A$ 是数域 $\mathbb K$ 上的 $n$ 阶方阵，求证：$A$ 的极小多项式的次数小于等于 $\operatorname{r}(A)+1$.

4. 设 $A$ 是数域 $\mathbb K$ 上的 $n$ 阶矩阵，求证：若 $\operatorname{tr}(A)=0$，则 $A$ 相似于一个 $\mathbb K$ 上的主对角元全为零的矩阵.

5. 设 $C$ 是数域 $\mathbb K$ 上的 $n$ 阶矩阵，求证：存在 $\mathbb K$ 上的 $n$ 阶矩阵 $A,B$，使 $AB-BA=C$ 的充分必要条件是 $\operatorname{tr}(C)=0$.

6. 设数域 $\mathbb K$ 上的 $n$ 阶矩阵 $A$ 的特征多项式 $f(\lambda)=P_1(\lambda)P_2(\lambda)\cdots P_k(\lambda)$，其中 $P_i(\lambda)$ ($1\leq i\leq k$) 是 $\mathbb K$ 上互异的首一不可约多项式. 设 $\mathbb K$ 上的 $n$ 阶矩阵 $B$ 满足 $AB=BA$，求证：存在 $\mathbb K$ 上次数不超过 $n-1$ 的多项式 $f(x)$，使 $B=f(A)$.

7. 设 $A$ 是数域 $\mathbb K$ 上的二阶矩阵，试求 $C(A)=\{X\in M_2(\mathbb K)\mid AX=XA\}$.

8. 设 $J=J_n(\lambda_0)$ 是特征值为 $\lambda_0$ 的 $n$ 阶 Jordan 块，求证：和 $J$ 乘法可交换的 $n$ 阶矩阵必可表示为 $J$ 的次数不超过 $n-1$ 的多项式.

9. 设数域 $\mathbb K$ 上的 $n$ 阶矩阵
\[
A=\left(\begin{array}{ccccc}
a_1 & b_1 & 0 & \cdots & 0\\
& \ddots & \ddots & \ddots & \vdots\\
& & \ddots & \ddots & 0\\
* & & & \ddots & b_{n-1}\\
& & & & a_n
\end{array}\right),
\]
其中 $b_1,\cdots,b_{n-1}$ 均不为零. 记 $C(A)=\{X\in M_n(\mathbb K)\mid AX=XA\}$，证明：线性空间 $C(A)$ 的一组基为 $\{I_n,A,\cdots,A^{n-1}\}$.

10. 设有 $n$ 阶分块对角阵
\[
A=\left(\begin{array}{ccc}
A_1 & & \\
& \ddots & \\
& & A_k
\end{array}\right),\quad
B=\left(\begin{array}{ccc}
B_1 & & \\
& \ddots & \\
& & B_k
\end{array}\right),
\]
其中 $A_i$ 和 $B_i$ 是同阶方阵. 设 $A_i$ 适合非零多项式 $g_i(x)$，且 $g_i(x)$ ($1\leq i\leq k$) 两两互素. 求证：若对每个 $i$，存在多项式 $f_i(x)$，使 $B_i=f_i(A_i)$，则必存在次数不超过 $n-1$ 的多项式 $f(x)$，使 $B=f(A)$.

11. 设 $n$ 阶矩阵 $A$ 的秩等于 $n-1$，$B$ 是同阶非零矩阵且 $AB=BA=O$，求证：存在次数不超过 $n-1$ 的多项式 $f(x)$，使 $B=f(A)$.

12. 设 $\varphi$ 是复线性空间 $V$ 上的线性变换，$V_0$ 是 $\varphi$ 的不变子空间. 求证：若 $\varphi$ 可对角化，则 $\varphi$ 在商空间 $V/V_0$ 上的诱导变换可对角化.

13. 设 $\varphi$ 是 $n$ 维复线性空间 $V$ 上的线性变换，求证：$\varphi$ 可对角化的充分必要条件是对任一 $\varphi$-不变子空间 $U$，均存在 $\varphi$-不变子空间 $W$，使得 $V=U\oplus W$. 这样的 $W$ 称为 $U$ 的 $\varphi$-不变补空间.

14. 设 $n$ 阶矩阵 $A$ 的极小多项式 $m(\lambda)$ 的次数为 $s$，$B=(b_{ij})$ 为 $s$ 阶矩阵，其中 $b_{ij}=\operatorname{tr}(A^{i+j-2})$（约定 $b_{11}=n$），求证：$A$ 可对角化的充分必要条件是 $B$ 为可逆阵.

15. 设 $n$ 阶复方阵 $A$ 的特征多项式为 $f(\lambda)$，复系数多项式 $g(\lambda)$ 满足 $(f(\lambda),g'(\lambda))=1$. 证明：$A$ 可对角化的充分必要条件是 $g(A)$ 可对角化.

16. 设 $V$ 为 $n$ 阶复矩阵全体构成的线性空间，$V$ 上的线性变换 $\varphi$ 定义为 $\varphi(X)=AXA$，其中 $A\in V$. 证明：$\varphi$ 可对角化的充分必要条件是 $A$ 可对角化.

17. 设 $V$ 为 $n$ 阶复矩阵全体构成的线性空间，$V$ 上的线性变换 $\varphi$ 定义为 $\varphi(X)=AX-XA$，其中 $A\in V$. 证明：$\varphi$ 可对角化的充分必要条件是 $A$ 可对角化.

18. 设 $\varphi$ 是 $n$ 维复线性空间 $V$ 上的线性变换，求证：$\varphi$ 可对角化的充分必要条件是对 $\varphi$ 的任一特征值 $\lambda_0$，总有 $\operatorname{Ker}(\varphi-\lambda_0I_V)\cap\operatorname{Im}(\varphi-\lambda_0I_V)=0$.

19. 求证：$n$ 阶复矩阵 $A$ 可对角化的充分必要条件是对 $A$ 的任一特征值 $\lambda_0$，$(\lambda_0I_n-A)^2$ 和 $\lambda_0I_n-A$ 的秩相同.

20. 若 $n$ ($n\geq 2$) 阶矩阵 $B$ 相似于 $R=\operatorname{diag}\left\{\left(\begin{array}{cc}0&1\\1&0\end{array}\right),I_{n-2}\right\}$，则称 $B$ 为反射阵. 证明：任一对合阵 $A$（即 $A^2=I_n$）均可分解为至多 $n$ 个两两乘法可交换的反射阵的乘积.

21. 设 $\varphi$ 是 $n$ 维线性空间 $V$ 上的线性变换，$U$ 是 $V$ 的非零 $\varphi$-不变子空间. 设 $\lambda_0$ 是限制变换 $\varphi|_U$ 的特征值，证明：$\varphi|_U$ 的属于特征值 $\lambda_0$ 的 Jordan 块的个数不超过 $\varphi$ 的属于特征值 $\lambda_0$ 的 Jordan 块的个数.

22. 设
\[
A=\left(\begin{array}{cccc}
1 & 0 & 0 & 0\\
a+2 & 1 & 0 & 0\\
5 & 3 & 1 & 0\\
7 & 6 & b+4 & 1
\end{array}\right),
\]
求 $A$ 的 Jordan 标准型.

23. 设 $J=J_n(0)$ 是特征值为零的 $n$ ($n\geq 2$) 阶 Jordan 块，求 $J^2$ 的 Jordan 标准型.

24. 求下列 $n$ ($n\geq 2$) 阶矩阵的 Jordan 标准型：
\[
A=\left(\begin{array}{cccccc}
c & 0 & 1 & 0 & \cdots & 0\\
& c & 0 & 1 & \cdots & 0\\
& & \ddots & \ddots & \ddots & \vdots\\
& & & \ddots & \ddots & 1\\
& & & & \ddots & 0\\
& & & & & c
\end{array}\right).
\]

25. 设 $\lambda_0$ 是 $n$ 阶矩阵 $A$ 的特征值，证明：对任意的正整数 $k$，特征值为 $\lambda_0$ 的 $k$ 阶 Jordan 块 $J_k(\lambda_0)$ 在 $A$ 的 Jordan 标准型 $J$ 中出现的个数为
\[
\operatorname{r}((A-\lambda_0I_n)^{k-1})+\operatorname{r}((A-\lambda_0I_n)^{k+1})-2\operatorname{r}((A-\lambda_0I_n)^k),
\]
其中约定 $\operatorname{r}((A-\lambda_0I_n)^0)=n$.

26. 设 $A,B$ 为 $n$ 阶矩阵，证明：它们相似的充分必要条件是对 $A$ 或 $B$ 的任一特征值 $\lambda_0$ 以及任意的 $1\leq k\leq n$，有
\[
\operatorname{r}((A-\lambda_0I_n)^k)=\operatorname{r}((B-\lambda_0I_n)^k).
\]

27. 设 $J=J_n(a)$ 是特征值为 $a\ne 0$ 的 $n$ 阶 Jordan 块，求 $J^m$ 的 Jordan 标准型，其中 $m$ 为非零整数.

28. 设 $J=J_n(0)$ 是特征值为零的 $n$ 阶 Jordan 块，求 $J^m$ ($m\geq 1$) 的 Jordan 标准型.

29. 设 $m$ 阶矩阵 $A$ 与 $n$ 阶矩阵 $B$ 没有公共的特征值，且 $A,B$ 的 Jordan 标准型分别为 $J_1,J_2$，又 $C$ 为 $m\times n$ 矩阵，求证：
\[
M=\left(\begin{array}{cc}
A & C\\
O & B
\end{array}\right)
\]
的 Jordan 标准型为 $\operatorname{diag}\{J_1,J_2\}$.

30. 设
\[
A=\left(\begin{array}{cccc}
1 & 0 & 0 & 0\\
b & a+1 & 0 & 0\\
3 & b & 2 & 0\\
5 & 4 & a & 2
\end{array}\right),
\]
求 $A$ 的 Jordan 标准型.

31. 设有 $n^2$ 个 $n$ 阶非零矩阵 $A_{ij}$ ($1\leq i,j\leq n$)，适合
\[
A_{ij}A_{jk}=A_{ik},\quad A_{ij}A_{lk}=O\quad (j\ne l).
\]
求证：存在可逆阵 $P$，使对任意的 $i,j$，$P^{-1}A_{ij}P=E_{ij}$，其中 $E_{ij}$ 是基础矩阵.

32. 设 $\lambda_0$ 是 $n$ 阶矩阵 $A$ 的特征值，其代数重数为 $m$. 设属于特征值 $\lambda_0$ 的最大 Jordan 块的阶数为 $k$，求证：
\[
\operatorname{r}(A-\lambda_0I_n)>\cdots>\operatorname{r}((A-\lambda_0I_n)^k)=\operatorname{r}((A-\lambda_0I_n)^{k+1})=\cdots=n-m.
\]

33. 设 $\lambda_0$ 是 $n$ 阶复矩阵 $A$ 的特征值，并且属于 $\lambda_0$ 的初等因子都是次数大于等于 2 的多项式. 求证：特征值 $\lambda_0$ 的任一特征向量 $\alpha$ 均可表示为 $A-\lambda_0I_n$ 的列向量的线性组合.

34. 设 $A,B$ 为 $n$ 阶矩阵，满足 $AB=BA=O$，$\operatorname{r}(A)=\operatorname{r}(A^2)$，求证：$\operatorname{r}(A+B)=\operatorname{r}(A)+\operatorname{r}(B)$.

35. 设 $A,B$ 分别是 $m,n$ 阶矩阵，求证：矩阵方程 $AX=XB$ 只有零解的充分必要条件是 $A,B$ 无公共的特征值.

36. 设 $A,B$ 分别是 $m,n$ 阶矩阵，$C$ 是 $m\times n$ 矩阵，求证：矩阵方程 $AX-XB=C$ 存在唯一解的充分必要条件是 $A,B$ 无公共的特征值.

37. 设 $A$ 为 $n$ 阶非异复矩阵，证明：对任一正整数 $m$，存在 $n$ 阶复矩阵 $B$，使 $A=B^m$.

38. 设 $A$ 为 $n$ 阶复矩阵，证明：存在 $n$ 阶非异复对称阵 $Q$，使 $Q^{-1}AQ=A'$.

39. 设 $A$ 为 $n$ 阶幂零阵，证明：$e^A$ 与 $I_n+A$ 相似.

40. 证明：存在 71 阶实方阵 $A$，使
\[
A^{70}+A^{69}+\cdots+A+I_{71}=\left(\begin{array}{ccccc}
2019 & 2018 & \cdots & 1949\\
& 2019 & \ddots & \vdots\\
& & \ddots & 2018\\
& & & 2019
\end{array}\right).
\]

41. 设 $a,b$ 都是实数，其中 $b\ne 0$，证明：对任意的正整数 $m$，存在 4 阶实方阵 $A$，使
\[
A^m=B=\left(\begin{array}{rrrr}
a & b & 2 & 0\\
-b & a & 2 & 0\\
0 & 0 & a & b\\
0 & 0 & -b & a
\end{array}\right).
\]

42. 设 $\varphi$ 为数域 $\mathbb K$ 上 $n$ 维线性空间 $V$ 上的线性变换，特征多项式与极小多项式分别为 $f(\lambda)$ 和 $m(\lambda)$，其不可约分解为：
\[
f(\lambda)=P_1(\lambda)^{r_1}P_2(\lambda)^{r_2}\cdots P_t(\lambda)^{r_t},\quad
m(\lambda)=P_1(\lambda)^{s_1}P_2(\lambda)^{s_2}\cdots P_t(\lambda)^{s_t},
\]
其中 $P_i(\lambda)$ 是 $\mathbb K$ 上互异的首一不可约多项式，$r_i>0$，$s_i>0$. 设 $V_i=\operatorname{Ker}P_i(\varphi)^{r_i}$，$U_i=\operatorname{Ker}P_i(\varphi)^{s_i}$，$1\leq i\leq t$. 求证：

(1) $V=V_1\oplus V_2\oplus\cdots\oplus V_t$，$U_i=V_i$ ($1\leq i\leq t$).

(2) $\varphi|_{V_i}$ 的特征多项式为 $P_i(\lambda)^{r_i}$，极小多项式为 $P_i(\lambda)^{s_i}$. 特别，$\dim V_i=r_i\deg P_i(\lambda)$.

43. 设 $n$ 维复线性空间 $V$ 上的线性变换 $\varphi$ 在一组基 $\{e_1,e_2,\cdots,e_n\}$ 下的表示矩阵为 Jordan 块 $J_n(\lambda_0)$，求所有的 $\varphi$-不变子空间.

44. 设 $\varphi$ 是 $n$ 维复线性空间 $V$ 上的线性变换，其特征多项式 $f(\lambda)$ 等于其极小多项式 $m(\lambda)$，求所有的 $\varphi$-不变子空间.

\end{document}
